\documentclass[11pt,a4paper]{article}

% ── Packages (only those available on this installation) ──────────────────────
\usepackage[margin=2.5cm]{geometry}
\usepackage{amsmath,amssymb}
\usepackage{graphicx}
\usepackage{booktabs}
\usepackage{longtable}
\usepackage{array}
\usepackage{xcolor}
\usepackage{hyperref}
\usepackage{caption}
\usepackage{subcaption}
\usepackage{enumitem}
\usepackage{microtype}
\usepackage{listings}
\usepackage{float}
\usepackage{placeins}

% Text extraction: map every glyph, ligatures included, to Unicode (review of 2026-09-28)
\input{glyphtounicode}
\pdfgentounicode=1

% ── Figure path ───────────────────────────────────────────────────────────────
\graphicspath{{../figures/}}

% ── Hyperref setup ────────────────────────────────────────────────────────────
\hypersetup{
  colorlinks = true,
  linkcolor  = blue!70!black,
  citecolor  = green!50!black,
  urlcolor   = blue!70!black,
}

% ── Custom commands ───────────────────────────────────────────────────────────
\newcommand{\numu}{$\nu_\mu$}
\newcommand{\numubar}{$\bar{\nu}_\mu$}
\newcommand{\ccpip}{$\mathrm{CC}\pi^{\pm}$}
\newcommand{\pmu}{$p_\mu$}
\newcommand{\ppi}{$p_\pi$}
\newcommand{\costhmu}{$\cos\theta_\mu$}
\newcommand{\costhpi}{$\cos\theta_\pi$}
\newcommand{\thmupi}{$\theta_{\mu\pi}$}
\newcommand{\uboone}{MicroBooNE}
\newcommand{\GeVc}{\,\text{GeV}\!/c}
\newcommand{\pot}{\,\text{POT}}
\newcommand{\cm}{\,\text{cm}}

% inline math shorthand so they work inside math mode too
\newcommand{\pmuM}{p_\mu}
\newcommand{\ppiM}{p_\pi}
\newcommand{\cthmuM}{\cos\theta_\mu}
\newcommand{\cthpiM}{\cos\theta_\pi}
\newcommand{\thmupM}{\theta_{\mu\pi}}

% ── Listings (C++ snippets) ───────────────────────────────────────────────────
\lstset{
  language=C++,
  basicstyle=\ttfamily\small,
  keywordstyle=\bfseries\color{blue!70!black},
  commentstyle=\itshape\color{gray},
  backgroundcolor=\color{gray!10},
  frame=single,
  framerule=0.4pt,
  numbers=none,
  breaklines=true,
}

% ── Table column helpers ──────────────────────────────────────────────────────
\newcolumntype{C}[1]{>{\centering\arraybackslash}p{#1}}
\newcolumntype{R}[1]{>{\raggedleft\arraybackslash}p{#1}}
\newcolumntype{L}[1]{>{\raggedright\arraybackslash}p{#1}}

% Let TeX stretch spacing on hard-to-break lines (long \texttt paths/code) to keep
% text out of the margin.
\emergencystretch=3em

% ─────────────────────────────────────────────────────────────────────────────

\usepackage{xr}
\externaldocument[N-]{analysis_note}
\externaldocument[S-]{technical_supplement}

\begin{document}

\begin{titlepage}
\centering
\vspace*{2cm}
{\LARGE\bfseries Change Log\par}
\vspace{0.5cm}
{\large Measurement of $\nu_\mu + \bar{\nu}_\mu$ Charged-Current Single Pion
Production on Argon in the NuMI Beam\par}
\vspace{0.8cm}
{\large Corrections, release repairs and documentation changes\par}
\vspace{0.5cm}
{\normalsize Draft 1.6, release 2026-09-19 \textperiodcentered\ document revision 2026-09-19\par}
\vspace{0.15cm}
{\footnotesize\ttfamily release 2026-09-06 \textperiodcentered\ 51 differential extractions $+$ 6 one-bin totals \textperiodcentered\ blind\par}
\vspace{1.5cm}
{\large Ishanee Pophale, Jarek Nowak\par}
\vspace{0.5cm}
{\normalsize\textit{Lancaster University}\par}
\vspace{2cm}

\begin{abstract}
This document records every correction that materially moved a number in the
$\mathrm{CC}\pi^{\pm}$ NuMI analysis, with the value before and after, together with the
release repairs and documentation changes made in the same period; entries of the latter
kind are labelled as such in their headings (``figure only'', ``presentation only'',
``documentation-only'').  It exists so
that a reader of the analysis note never has to work out whether a given table predates
a given fix, and so that a result quoted from an older draft can be identified as
superseded rather than merely different.

Most entries are normalisation errors. The extraction divides by four numbers (the data
exposure, the integrated flux, the argon target count and the per-run Monte Carlo exposure), and
eight separate defects came from one of those four differing between configurations without
anything detecting it.
The response is \texttt{norm\_manifest.py}, which reads the authoritative configuration
files and asserts that they agree, so that a future drift fails a check instead of
producing a wrong cross section.

Three corrections changed physics conclusions rather than only normalisation, and the
conclusions they invalidated are marked as withdrawn here and in the analysis note: the
non-uniform-grid regularisation defect, which had collapsed several additional-smearing
matrices towards rank one; the beam-axis convention, under which reconstructed angular and
transverse observables were not the same quantities as the generator predictions they were
being compared against; and the double application of $A_C$ to the model predictions,
which invented a muon-angle anomaly.
\end{abstract}
\end{titlepage}

\tableofcontents
\clearpage

\section{How to read this document}

Each entry states what was wrong, how it was found, what changed numerically, and what, if
anything, had to be withdrawn.  ``Before'' values are quoted from the draft in
which they appeared and should not be used; they are given so that an older table can be
recognised.

Entries are ordered by the stage of the analysis they affect, not chronologically:
exposure and flux first, then the Monte Carlo weighting, then the unfolding, then the
systematic covariance.

\section{Exposure, flux, and target normalisation}

\subsection{Obsolete integrated flux (factor ${\sim}3.48$)}

An outdated flux file was used to convert the extracted event rate into a cross section.
The integrated flux it supplied was larger than the correct value by a factor of
approximately $3.48$, so every cross section produced with it was smaller than the truth
by the same factor.

\textbf{Effect.} All cross sections, all configurations.  Any absolute normalisation from
a draft predating this fix is wrong by ${\sim}3.48$ and must not be quoted.  Shapes were
unaffected, because the error was a single multiplicative constant.

\subsection{Configuration-specific flux was not being used}

The RHC and combined extractions were normalised with the FHC integrated flux.  The three
exposures see different spectra, and the correct integrated fluxes are
\[
\Phi_{\mathrm{FHC}} = 6.81159\times10^{-10},\quad
\Phi_{\mathrm{RHC}} = 6.44646\times10^{-10},\quad
\Phi_{\mathrm{comb}} = 6.60865\times10^{-10}
\]
in $\mathrm{cm}^{-2}\,\mathrm{POT}^{-1}$.

\textbf{Effect.} RHC cross sections rose by $5.7\%$ and combined by $3.1\%$.  FHC was
already correct and did not move.  The flux is now set by an explicit \texttt{Flux} line
in each cross-section configuration file rather than inherited, and
\texttt{norm\_manifest.py} checks that all configurations of a given exposure agree on it.

\subsection{Per-mode rather than per-run exposure weighting}

Data and Monte Carlo exposures were pooled per horn mode and a single scale factor
applied.  Because the data-to-Monte-Carlo exposure ratio differs substantially between
runs, this weighted the result by how much simulation happened to exist for each run
rather than by the data actually collected.  Each run is now scaled by its own
$S_r = \mathrm{POT}^{\mathrm{data}}_r/\mathrm{POT}^{\mathrm{MC}}_r$.

\textbf{Effect.} Changed the relative weight of runs within every configuration.  This is
also the change that makes the framework ready for real data, where the per-run exposures
are not free parameters.

\subsection{Duplicated run-total POT in split files}

Two conventions coexist in the processed samples.  In \texttt{processed/} a run split
across several files repeats the \emph{run-total} exposure in every file; in
\texttt{processed/w/} each file carries its own \emph{true} exposure.  Dividing by the
per-file value broke the first convention and summing over files broke the second: the
first made the inclusive tune come out a factor of two low, the second drove an RHC
proton-tagged cross section negative.

The resolution is to normalise by the sum of \emph{distinct} per-run exposures, which is
correct under both conventions.

\textbf{Effect.} At the time of the fix the RHC $W_{\pi p}$ cross section moved from
$-0.194$ to $+0.578$ (in units of $10^{-38}\,\mathrm{cm}^2/\mathrm{Ar}$), and the
inclusive RHC $p_\mu$ result was unchanged at $0.860$.  Results retracted on the strength
of the negative value were reinstated.  Both numbers have since moved again under the
central-value weighting fix below; the current values are in the analysis note.

\section{Monte Carlo weighting}

\subsection{Central-value weight applied twice to the fake data}

This is the largest single correction and it affected every result.

The fake data are generated externally by \texttt{throw\_perrun\_*.C}, which folds the
tuned central value, the PPFX central value and the normalisation into the event
multiplicity and then writes the per-event weights back as unity: the throw is a
Poisson realisation of an already-weighted expectation.  The framework's
\texttt{build\_universes} nevertheless substituted the \emph{weighted} central-value
histogram for the raw counts, on both the reconstructed and the truth side, so the
central value was applied a second time.

\textbf{How it was found.} By asking directly whether the unfolded result was being
compared against the input that generated it.  Four independent checks confirmed the
diagnosis: $1621.0$ is a direct count of selected events across the four FHC files;
$1721.9$ equals both \texttt{FakeDataMC\_0\_reco} and
$\texttt{onBNB\_reco}-\texttt{extBNB\_reco}$; run~1 gives $558.0$ entries unweighted
against $599.9$ weighted for the same $558$ events; and the throw's own expectation of
$1567.1$ agrees with the thrown count at $1.3\sigma$.

\textbf{Effect.} The mean inclusive closure ratio moved from $1.106$ to $1.001$ ($1.016$ once the truth-side correction below is included, $1.018$ on the re-thrown fake data), and all
$51$ extracted cross sections fell by between $3.8\%$ and $12.9\%$.  Every absolute cross
section from a draft predating this fix is high by roughly a tenth.

\textbf{Withdrawn.} Two earlier explanations of the closure excess, one attributing it
to the unfolding covariance and Wiener filter and one to $A_C$ smearing, were both wrong
and are retracted.  They are retained in the technical supplement as recorded dead ends
rather than deleted.

\subsection{Generator prediction histograms truncated at the overflow}

Generator predictions were read without their overflow contents, so the highest-momentum
region was silently missing from the comparison curves.

\textbf{Effect.} Model comparison only; the measured cross sections were unaffected.  Any
generator $\chi^2$ from a draft predating the regeneration is invalid.

\section{Unfolding}

\subsection{Second-derivative regularisation on a non-uniform grid}

The Wiener-SVD second-derivative regulariser used the uniform-grid $(1,-2,1)$ stencil on
bins whose widths differ by factors of ten to fifteen.  On such a grid that operator is
not an approximation to the second derivative of anything, and its near-null direction is
not the physically flat one.  The Wiener filter can then retain that direction alone,
which drives the additional-smearing matrix $A_C$ towards rank one and destroys the shape
information the measurement is supposed to carry.

The corrected operator differences the \emph{density} on a non-uniform grid, reducing to
the old stencil when the widths are equal, with one-sided first-derivative rows at the
ends so that the operator stays non-singular without asserting that the distribution
vanishes at the edges, which it does not, since $\cos\theta_\mu$ peaks at the upper
edge.

\textbf{Effect} on the conditioning measure $s_2/s_1$ of $A_C$:
\begin{center}
\begin{tabular}{lcc}
\toprule
Observable & before & after \\
\midrule
FHC $\theta_\mu$              & $0.015$ & $0.51$ \\
combined $\theta_\mu$         & $0.045$ & $0.38$ \\
proton-tagged $W_{\pi p}$     & $0.022$ & $0.60$ \\
proton-tagged $p_n$           & $0.21$  & $0.88$ \\
\bottomrule
\end{tabular}
\end{center}

\textbf{Withdrawn.} The earlier conclusion that $\theta_\mu$ ``repairs'' what
$\cos\theta_\mu$ cannot is retracted: under the defective operator every model was being
collapsed onto a common shape, so the apparent agreement was an artefact of the
regularisation rather than a property of the data.  After the fix $\theta_\mu$
discriminates between models again and prefers the tune that generated the fake data.

This entry is the clearest illustration of why closure alone is not a sufficient
validation.  A rank-one extraction closes perfectly against its own correspondingly
smeared truth while retaining almost no shape information.

\subsection{Angular and transverse observables were defined about detector \texorpdfstring{$z$}{z}}

The NuMI neutrinos arrive approximately $28.6^\circ$ away from the detector $z$ axis, but
reconstructed angles and transverse kinematic imbalance were being computed about $z$
while the generator predictions were computed about the neutrino direction.  These are
not the same observables.  Truth-level quantities now use the per-event neutrino
direction and reconstruction uses the target-to-vertex direction; the two constructions
differ only at the milliradian level.

\textbf{Effect.} $64.6\%$ of signal events change $\cos\theta_\mu$ bin and $44.6\%$ change
$\cos\theta_\pi$ bin.  The mean $\delta p_T$ moves from $0.742$ to
$0.264\ \mathrm{GeV}/c$.  The muon and pion momenta, the opening angle $\theta_{\mu\pi}$,
the hadronic mass $W_{\pi p}$, and the event selection are all unaffected.

\textbf{Withdrawn.} All proton-tagged transverse-kinematic model comparisons made before
this correction are withdrawn, because the measured and predicted quantities were
different observables.

\subsection{Predictions were smeared by \texorpdfstring{$A_C$}{A C} twice}

Every model prediction plotted or compared against the data was multiplied by the
additional-smearing matrix twice: once by the extractor, which applies $A_C$ to each entry
of the prediction map in place, and again on the plotting path.  Confirmed exactly, as
$A_C^{2}\cdot\text{raw} = \text{plotted}\times\text{width}$ to $1.0000$ in every bin.

\textbf{How it was found.}  By testing whether the released $A_C$ lets an external user
reproduce a published generator curve. It did not: the discrepancy was $31\%$ on the
integral and varied from $1.21$ to $1.55$ across bins, so it was not a normalisation, and
the residual factor turned out to be $A_C$ itself.

\textbf{Effect} (2026-08-31 pass; the current values are in the note's
$\chi^2$ tables).  Model $\chi^2$ values only.  The anomaly the note previously attributed
to $A_C$ suppression of sharply peaked generators was this bug:

\begin{center}
\small
\begin{tabular}{llcc}
\toprule
Observable & Model & before & after \\
\midrule
FHC $\theta_\mu$      & NEUT  & $40.60/5$ & $8.96/5$ \\
FHC $\theta_\mu$      & NuWro & $38.62/5$ & $9.43/5$ \\
FHC $\cos\theta_\mu$  & NEUT  & $29.78/5$ & $4.90/5$ \\
FHC $\cos\theta_\mu$  & NuWro & $29.10/5$ & $4.93/5$ \\
FHC $p_\mu$           & GENIE & $5.52/7$  & $0.57/7$ \\
\bottomrule
\end{tabular}
\end{center}

\textbf{Not affected.}  The unfolded cross sections (FHC $p_\mu$ integrates to $0.8275$
before and after) and the closure test, because the fake-data truth was the one curve not
passed through the second multiplication.  The previous table also predated the
central-value weighting fix, so part of the shift in its ``truth'' column is that
correction rather than this one.

\textbf{Withdrawn.}  The argument for preferring $\theta$ over $\cos\theta$ on the grounds
that the $\cos\theta_\mu$ $\chi^2$ was an artefact.  The effect that motivated it does not
exist.

\section{Systematic covariance}

\subsection{Combined detector-variation exposure double-counted}

The combined configuration pooled detector-variation samples per horn mode, and the
denominator used a single file's exposure rather than the summed Monte Carlo exposure over
all files of that variation type, so the combined detector systematic was inflated.

\textbf{Effect.} The total detector systematic fell from $42.3\%$ to $24.6\%$, and the
combined result moved from $0.910$ to $1.027$ times its expected value.  FHC and RHC were
bit-identical before and after, which is the check that isolates the defect to the
combination step.

\subsection{Proton-tagged detector systematics were absent entirely}

Of the fifty-one extractions, the thirty-three proton-tagged ones carried \emph{no}
detector systematic at all: the \texttt{detVar} column was \texttt{n/a} for every
proton-tagged observable in every configuration, so their quoted uncertainties omitted a
term worth $17$--$43\%$ in the corresponding inclusive results.

The cause was a configuration gap rather than missing data.  The proton-tagged detector
variation samples had been processed on 2026-08-20/21 and carry the required branches, but
nothing referenced them: the \texttt{\_w} file-properties tables contained no
\texttt{detVar} lines, and \texttt{ccpi1p\_systcalc\_numi.conf} had no \texttt{DV} block
at all.  Both were needed; adding either alone does nothing.

\textbf{Effect} (values of the time).  The thirty-three proton-tagged universes were rebuilt
with the nine detector variation samples included and re-unfolded from them.  Every
extraction in the analysis now carries a detector term.  The proton-tagged terms ran
$10.5\%$--$64.4\%$ with a median of $23.9\%$, against $16.9\%$--$42.8\%$ for the inclusive
results (release of 2026-09-26: $9.9$--$74.5\%$, median $30.9\%$, and $16.6$--$34.2\%$).

\textbf{Check} (of the time).  Per-mode detector-variation double counting has occurred in
this analysis before, and the combined configuration, which carries two per-mode sets
where the single-mode configurations carry one, is the one exposed to it.  Double
counting inflates the combined result.  On that pass, of the eight observables (of eleven)
whose combined value fell outside its own FHC--RHC range, six fell \emph{below} both and
only two above, the finite-sample averaging-down expected of statistically limited
detector-variation samples rather than a recurrence; on the current release only two of
eleven fall outside the range, both above (analysis note).

\section{Count-level normalisation, fake data and release repairs (2026-09-02 onwards)}

\subsection{Cut-and-count uncertainty quoted from a pre-rebuild breakdown}

The total flux-averaged cross sections of the cut-and-count path took their systematic
fraction from the $p_\mu$ systematic-breakdown tables, which had not been regenerated
after the detector-variation rebuild of 2026-08-31 (their detector terms were
$\sim10\%$ against $17$--$43\%$ in the release).  The breakdown tables in the note and the
supplement were regenerated from the release systematics dumps on 2026-09-02 and the
cut-and-count uncertainties recomputed from them; the macro also now applies the
central-value MC weight, which moved the central values by $0.3$--$0.5\%$.

\textbf{Effect} (values of the time; the current totals are in the one-bin entry below).  Inclusive totals ($10^{-38}$~cm$^2$/Ar): FHC $1.06\pm0.27\to1.03\pm0.43$,
RHC $0.98\pm0.35\to0.93\pm0.42$, combined $1.02\pm0.30\to0.98\pm0.40$.  The generator
comparison changes from ``GENIE $\sim1\sigma$ low'' to ``GENIE $0.5$--$0.6\sigma$ low''.
A proton-tagged cut-and-count total is quoted for the first time (FHC $0.63\pm0.31$,
RHC $0.55\pm0.21$, combined $0.59\pm0.25$).

\subsection{Fake-data truth still carried the central-value weight}

The 2026-08-30 correction (above) stopped the tune weight being applied a second time
to the externally thrown fake data on the reconstructed side.  The truth side was
missed: the throw macro resets only the scalar weight branches, the
\texttt{weight\_TunedCentralValue\_UBGenie} vector survives, so the universe maker still
writes a tune-weighted ``CV'' universe for each fake-data file, and the systematics
calculator took that weighted histogram as the fake-data \emph{truth} whenever it was
present.  Measured on the FHC $p_\mu$ universe file: weighted/unweighted true-signal
integrals $1.023$, $1.008$, $1.008$, $1.018$ for Runs~1, 2, 4, 5, about $1.6\%$ overall.
Every closure ratio was therefore low by that amount and every closure $\chi^2$ slightly
too small.  Found in the 2026-09-02 code trace; the truth path now uses the same
external-throw flag as the reco path, and the universe cache key was bumped so the
stale sums cannot be reused.

\textbf{Effect.}  Unfolded cross sections and covariances are unchanged (FHC $p_\mu$
integrates to $0.8275$ before and after).  Closure: FHC $p_\mu$ unfolded/truth
$1.024\to1.041$, $\chi^2$ $0.10/7\to0.13/7$; the mean inclusive closure ratio
$1.001\to1.016$.  All fifty-one sidecars, the released curves
(\texttt{truth\_smeared} column), the closure tables and the closure figures were
regenerated; the ensemble bias is unaffected because its reference is the fixed
central-value truth, not the thrown one.

\subsection{Beam-off gate count was an event count}

Every live file-properties configuration scaled the pooled Run-1$+$Run-3b beam-off
sample by $\text{bnb\_trigs}[\text{run}]/3\,821\,593$.  The denominator was inherited from
the earlier custom selection, where it was the sum $610\,496+3\,211\,097$ of the two
samples' \emph{event} counts; their gate counts, from the analyser's trigger table
supplied on 2026-09-02, are $4\,582\,248.27+32\,649\,128.65=37\,231\,376.92$, a factor
$9.74$ larger.  The selected cosmic rate per gate exposed it: $5.8$ per $10^{6}$ gates for
the pooled sample against $0.37$--$1.3$ for the dedicated per-period files, versus $0.59$
with the gate count.  Corrected in all six live configurations and in the counting and
cut-flow macros; the universes were re-summed (the cache key tracks the file-properties
contents) and all fifty-one extractions, covariances and the release were regenerated.

\textbf{Effect.}  EXT background at the FHC final cut $128\to13$ events (RHC
$133\to14$, combined $261\to27$); purity $55.8\to59.7\%$ (FHC).  Fake-data central values
are unchanged by construction (the same EXT is added and subtracted); the EXT-statistics
and data-statistics covariance terms shrink, which moves the Wiener-SVD results only
through the regularisation (see the regenerated tables).  A real-data extraction with the
old value would have over-subtracted $\approx\!115$ events, $\approx\!12\%$ of the FHC
signal.  \textbf{Withdrawn:} every beam-off yield, purity and EXT-statistics term quoted
in drafts before 1.2.  The sideband tables were regenerated with the per-run treatment
(\texttt{regen\_sidebands.sh}); the cosmic control region, formerly EXT-dominated by
$4.5:1$, now has $0.7$ EXT events per neutrino-MC event.

\subsection{Beam-off matched by run period instead of one pooled sample}

Beam-off data are cosmic-only, so a sample is matched to a beam-on run by its run period,
irrespective of the horn polarity it was recorded under (analyser, 2026-09-02).  The
pooled Run-1$+$Run-3b sample was replaced by the dedicated per-period samples (Run~1;
Run~3b; Run~4a, 4b, 4c, 4d pooled; Run~5; the Run-1 sample standing in for Run~2), each
scaled by its own gate count from the analyser's table, and processed with the swtrig
fail-open build so that the CRT-era samples select cosmics at all.  The selected rate per
$10^{6}$ gates is $1.75$ (Run~1), $0.43$ (Run~3b), $1.15$ (Run~4), $0.52$ (Run~5), so
period matching moves the prediction by a factor of about two relative to the pooled
sample.  All fifty-one universe files were rebuilt (SLURM array 3205007, 51/51), the
extractions re-unfolded, covariances re-harvested and the release regenerated.

\textbf{Effect.}  EXT background at the final cut (FHC / RHC / combined): pooled
$13/14/27$ $\to$ per-run $30/23/53$ events, i.e.\ $1.9\%$ of the selected sample; purity
$59.1\%$ (FHC).  EXT-statistics term on FHC $p_\mu$ $0.4\to0.9\%$.  Unfolded central
values move by $\lesssim1\%$ through the regularisation.

\subsection{PPFX central-value weight absent from the cut-flow and cut-and-count counts}

The framework's central-value event weight for NuMI is tune $\times$ PPFX~CV $\times$
normalisation (\texttt{UniverseMaker.cxx}), and the per-run fake data are thrown with
that product.  The selection's cut-flow histograms and the cut-and-count macro carried
the tune only.  The PPFX central-value weight averages $0.944$ on selected signal, so
every absolute yield in the cut-flow table and every cut-and-count total sat $3$--$6\%$
above the prediction the unfolding uses; efficiencies and purities, being ratios, were
unaffected at the $0.5\%$ level.  Found in the 2026-09-02 code trace.

\textbf{Effect.}  Cut-and-count totals: FHC $1.06\to1.03$, RHC $0.98\to0.93$, combined
$1.02\to0.98$ ($10^{-38}$~cm$^2$/Ar), and the RHC/FHC ratio $0.928\to0.901$.  The
selection now fills the cut-flow histograms with the full weight and the fourteen
overlay files and the dirt overlay were reprocessed on 2026-09-03
(\texttt{reprocess\_cutflow\_cf.sh}, into a separate directory so the universes' POT
convention is untouched); the cut-flow table is regenerated from them.  No unfolded
result is affected.

\subsection{Proton-tagged TKI figure drawn from the withdrawn fine-binned sidecars}

The proton-tagged W/TKI montage macros read closure sidecars keyed \texttt{dpt},
\texttt{dalphat}, \texttt{dphit}, \texttt{pn}: the withdrawn fine-binned extraction of
2026-08-18, which predates the beam-axis correction, rather than the released two-bin
\texttt{*2bin} files.  The figure therefore compared transverse quantities about detector
$z$ (its ``data'' and tune) with beam-frame generator curves, and showed shape
disagreements of a factor of several that do not exist in the release.  The staleness
check maps figures by configuration only and could not catch it.  Macros repointed, the
stale sidecars quarantined, the figure regenerated (2026-09-03).

\textbf{Effect.}  Figure only; no released number.  In the two-bin scheme the generators
differ by normalisation ($\times1.6$ GENIE to NEUT) with shapes consistent with the data.

\subsection{Result tables rewritten into the wrong place by the table patcher}

The script that repatched the result tables after each regeneration located a table by
the \texttt{\textbackslash begin\{tabular\}} nearest its label, falling back to the
\emph{preceding} tabular whenever the caption between label and tabular exceeded 400
characters.  Four tables were affected in the committed Draft~1.2: the FHC results table
received the RHC rows, the RHC table the combined rows, the combined table the rows of the
integrated-cross-section summary, and in the supplement the cut-and-count comparison was
overwritten by an FHC systematics breakdown.  Caught by the second external review
(2026-09-03).  All release-derived tables are now regenerated as whole tabulars by
\texttt{rebuild\_tables.py} with a locator that requires the tabular to follow the label
inside the same environment, and \texttt{report/check\_tables.py} asserts, on every build,
that each table sits under its own label, has the column count its header declares, and
agrees with \texttt{current\_results.tsv}.

\textbf{Effect.}  Presentation only; the release files were correct throughout.  The
tables of the 2026-09-03 build are the first to pass the automated check.

\subsection{Two central-value weight rules were in use; the fake data skipped events the prediction counted}

The framework's \texttt{safe\_weight()} sets a non-finite, negative or $>30$ central-value
weight to unity. The count-level macros (cut-and-count, cut-flow, sidebands, transfer
factors) instead zeroed invalid weights and did not clamp, and the per-run fake-data throws
\emph{skipped} events with a non-finite weight. About $1.15\%$ of overlay events carry a
non-finite tune weight, concentrated in the non-signal $\nu_\mu$CC classes, so the
pseudo-data sat below the framework's own prediction by $0.9\%$ in the signal region
($13$--$14$ events) and by $2.8$--$3.4\%$ in the CC$0\pi$ control region ($341$ FHC,
$422$ RHC POT-scaled events); the sideband yields of the note differed from the
systematics tables by the same amounts. Found through the third external review
(2026-09-03).

\textbf{Correction.} One rule, the safe-weight definition in the Monte Carlo section of the note, is now applied
everywhere: the selection's cut-flow histograms (reprocessed), every count-level macro,
and the fake-data throws, which were re-thrown with the same seeds and the events kept at
unit weight; all fifty-one universe files were rebuilt from the re-thrown fake data.
\textbf{Effect.} Pseudo-data yields rise by $0.9\%$ (signal region) to $3.4\%$ (CC$0\pi$).
Because the re-throw is a new Poisson realisation, the individual Wiener-SVD integrals
moved by up to $20\%$ (RHC $p_\mu$ $0.573\to0.690$) within their statistical
uncertainties, the inclusive closure mean moved from $1.016$ to $1.018$ (all $p>0.9$), and
the prediction-total uncertainty that the cut-and-count totals inherit from the $p_\mu$
extraction moved from $44.4\%$ to $37.7\%$ (RHC) and from $37.0\%$ to $53.5\%$
(proton-tagged RHC); the cut-and-count central values, which use the central-value
prediction rather than the throw, are unchanged. The D'Agostini cross-check was re-run on
the new release: its integral is $20$--$29\%$ above Wiener-SVD and flat across observables
to $\le1.0\%$ (previously $27$--$33\%$ and $\le2.1\%$). The real-data path was never affected.

\subsection{The total cross section borrowed its systematic from the $p_\mu$ extraction}

The headline total was a cut-and-count evaluation whose systematic fraction was taken from
the prediction total of the released $p_\mu$ extraction of the same configuration. That
proxy carried the regularisation amplification of the flux and detector terms, and it
depended on the fake-data throw through the unfolding operator: on the re-throw of
2026-09-03 alone it moved from $44.4\%$ to $37.7\%$ (RHC) and from $37.0\%$ to $53.5\%$
(proton-tagged RHC) with no change to the central value.

\textbf{Correction.} The total is now extracted with the framework as a one-bin
distribution: one true bin containing every signal event, defined on $\cos\theta_\mu$
over its full range so that no overflow exists, one reconstructed bin, Wiener filter off.
The response is the efficiency alone, nothing is regularised, $A_C\equiv1$ exactly, and
the covariance is propagated universe by universe like every other extraction. Six new
universe files (inclusive and proton-tagged $\times$ FHC/RHC/combined; SLURM 3228639) and
\texttt{data\_release/total\_xsec.tsv}. The cut-and-count macro is kept as the
closed-form check and reproduces the one-bin extraction of the central-value prediction to
four digits.
\textbf{Effect.} Inclusive totals $1.03\pm0.40$, $0.93\pm0.35$, $0.97\pm0.37$
(previously $\pm0.42$, $\pm0.35$, $\pm0.34$ with the borrowed systematic); proton-tagged
$0.62\pm0.26$, $0.53\pm0.25$, $0.57\pm0.24$ (previously $0.63\pm0.32$, $0.55\pm0.30$,
$0.59\pm0.28$). The central values now fluctuate with the throw (closure $0.98$--$1.00$
against the realised truth) instead of being the injected tune by construction. The
generator comparison moves from ``$0.5$--$0.6\sigma$ above GENIE, within $0.7\sigma$ of
the rest'' to ``$0.5$--$0.8\sigma$ above GENIE, within $0.8\sigma$ of the rest''.

\section{Development-era corrections}

These predate the framework-based analysis and are recorded for completeness.  They were
previously listed in a ``Known Issues'' table in the analysis note; that table had become
thirteen resolved rows against one open one, so it was retired and its contents moved
here, where a resolved correction belongs.

\begin{center}
\small
\begin{tabular}{p{0.36\textwidth}p{0.56\textwidth}}
\toprule
Item & Resolution \\
\midrule
Unfolding covariance was statistics-only
  & Replaced by the framework \texttt{PredTotal} plus the data Poisson term, removing a
    ${\sim}26\%$ bias in the quoted uncertainty. \\
$\theta_{\mu\pi}$ binning mismatch
  & A custom 10-bin scheme was aligned to the framework's 9-bin one; the two pipelines had
    been binning the same observable differently. \\
Negative tail bins
  & An artefact of the EXT treatment; all bins positive once the subtraction was
    corrected. \\
EXT subtraction on fake data
  & The custom pipeline cancelled it via \texttt{fake\_data\_ext\_cancel}, a flag that had
    to be remembered at unblinding. Superseded: the framework selects the beam-off
    treatment from the presence of Monte Carlo truth in the beam-on file (added to
    pseudo-data, subtracted from real data) and refuses a real file unless
    \texttt{XSEC\_UNBLIND=1}, so no manual switch remains. \\
Multisim propagation in the custom pipeline
  & Never completed, and now moot: the released results come from the framework, which
    propagates the multisims (\texttt{cov\_xsec\_multi} is in the release).  The row
    survived in the note as the single ``open'' item long after it had ceased to describe
    anything. \\
CRT veto branches
  & \texttt{crthitpe} and \texttt{crtveto} are all-zero in the NuMI ntuples; replaced by
    \texttt{bdt\_cosmic}, disabled at its default. \\
\bottomrule
\end{tabular}
\end{center}

\section{Review passes}

\subsection{Corrections from the independent review of 2026-09-04}

Two reviewers with no prior contact with the analysis read Draft~1.4 in full, one checking
every number against the release files and one as a referee. What they found, and what
changed:

\textbf{Release index.} \texttt{data\_release/index\_extractions.tsv} and \texttt{README.md}
were the 2026-08-31 files and disagreed with every other release artefact; regenerated
from the current results (\texttt{report/tools/index\_extractions.py}), with the
differential $W_{\pi p}$ shape flagged withdrawn.

\textbf{Proton-tagged cut-flow.} The supplement's proton-tagged cut-flow table predated
the beam-off gate and PPFX-weight corrections, and the note's proton-tagged prose quoted
it ($47$--$48\%$ purity, $33\to48\%$). The fourteen overlays and the dirt overlay were
reprocessed with the proton-tagged selection under the unified weight rule
(\texttt{processed/cf\_1p}); the macro's final stage, which had counted events
unweighted ($2.6\%$ high), now carries the same weight as the other stages, and its
final row reproduces the cut-and-count check exactly ($376.4$ signal, $723.0$ predicted,
FHC). Purity is $52\%$ (FHC), $50\%$ (RHC), $51\%$ (combined); the tag rejects $67\%$ of
the neutrino background for a $32\%$ signal loss and reduces the beam-off by
$3.3$--$3.5$ at the final cut, not six.

\textbf{Flux term.} The note attributed the $24$--$33\%$ unfolded flux term to amplification
by the unfolding operator. The one-bin total, which has no regularisation, carries
$28$--$30\%$: a flux universe scales signal and background together, so the term is
$\delta\Phi/\Phi\times(1+B/S)$, $17.7\%\times1.71=30.2\%$ against the one-bin $30.25\%$
(FHC). The section is rewritten; the reconstructed-level table of the earlier pass is
dropped.

\textbf{Control-region procedure.} The target class of the $\pi^0$ region as tabulated
($\nu_\mu$CC with a $\pi^0$) is $46\%$ of the region, below the $50\%$ purity the
procedure requires for a fitted correction; stated, with the CC$+$NC~$\pi^0$ definition
to be ratified. The pass/fail criterion is now a global normalisation $\chi^2$ over the
eight (region, mode) totals; the sixteen independent $2\sigma$ and $p>0.05$ tests of the
earlier draft (which would have flagged at least one region by chance more often than
not) are diagnostics, reported with pulls against the statistical$+$detector covariance
and a normalisation-marginalised shape test. The response to a localised reconstruction
failure, the joint FHC/RHC treatment and the leverage of the held-fixed signal
contamination are specified. A control-region skim that never writes signal-region
events is added to the prerequisites; it did not exist at that review and was implemented
on 2026-09-06 (below).

\textbf{Adopted 2026-09-06, pending the conveners' ratification.} The analysers adopted the $\pi^0$ region's target class as
CC$+$NC~$\pi^0$ (purity $50.1\%$ FHC, $51.0\%$ RHC of the neutrino MC; $T=0.019$--$0.020$;
transfer uncertainty $17$--$19\%$, floor-limited); the class was added to the transfer,
transfer-uncertainty and detector-variation macros and to the note's tables, and the
purity denominator of the procedure is now defined as the neutrino-MC prediction.

\textbf{Protocol statistics (2026-09-06).} The three statistics of the frozen procedure
(global $8\times8$ normalisation $\chi^2$, per-region pulls against the full and the
statistical$+$detector covariance, normalisation-marginalised shape tests) were implemented
(\texttt{macros/sb\_protocol.C}) and exercised on the full-stack pseudo-data: $\chi^2$
$0.66/8$, pulls within $\pm0.5$, shape $p>0.47$; power $12\%$ (CC$0\pi$, $\pi^0$) and
$50$--$62\%$ (multi-$\pi$, cosmic) for a single-region shift, $58\%$ for a common shift.
The response tree was rewritten as an explicit flow (trigger, order of regions, numeric
criterion for a localised reconstruction failure, size of the added uncertainty, leverage
of the held-fixed components in a fitted correction, post-treatment criterion).

\textbf{Detector term statistics.} The Monte-Carlo-statistical floor of the
detector-variation term was measured on split samples
(\texttt{macros/sb\_detvar\_split.C}): the CC$0\pi$ transfer term ($27\%$ FHC, $31\%$
RHC) is genuine (floor $9\%$, $6\%$); the $\pi^0$ transfer term ($15\%$, $11\%$) is at or
below its floor ($13\%$, $17\%$) and is an upper limit; the region-yield terms have floors
below $2\%$. The procedure now uses the larger of the term and its floor.

\textbf{Ensemble.} The 32-throw ensemble was thrown by a macro that still skipped
non-finite-weight events, which accounts for $\approx1.4\%$ of its $2.5\%$ negative offset.
The macro now applies the unified rule and the ensemble was re-run for $p_\mu$,
$\cos\theta_\mu$ and the one-bin total (SLURM 3277674, 3277691, 3277692; statistical
covariances re-harvested with \texttt{slurm/ens\_statcov.sh}, pulls with
\texttt{report/tools/ensemble\_stat\_pulls.py}). Statistics-only pull widths $1.08$,
$0.97$, $1.02$ (previously $1.04$, $0.96$); ensemble means below the reference by $0.5\%$,
$2.7\%$ and $1.0\%$ at $0.4$, $1.9$ and $1.3\sigma$ (previously $2.5\%$ at $2.1\sigma$
for $p_\mu$). The bias statement in the note is reworded accordingly: a small common
negative offset, not significant, largest in the most regularised observable, and present
at the $1\%$ level in the unregularised total.

\textbf{Also.} The $A_C$ claims made about proton-tagged closure ratios (which are ratios to
the $A_C$-smeared truth) are removed; the RHC-below-FHC comparison of $A_C$-suppressed
integrals is replaced by the one-bin ratio; the reproducibility claim is scoped to the
2026-09-02 pass; the D'Agostini cross-check is stated to be prior-dominated on
central-value input; the absence of an alternative-generator fake-data test on the
released chain, of an MCS momentum systematic and of gate-ratio and Run-2 stand-in
beam-off terms is stated in the Limitations; the signal definition (one charged pion of
any momentum, above threshold), the unfolder provenance, the regularisation matrix and
the proton-tag values ($\mathrm{LLR}<0.05$, $p_p>0.3$~GeV/$c$) are stated precisely; the
flux mean energies ($0.48$~GeV) and the fraction above $0.3$~GeV ($40\%$) are given, which
is the origin of the factor $\sim\!3$ against the BNB total. Some twenty further inline
numbers that had survived from older passes were corrected.

\subsection{Inclusive $\cos\theta_\pi$ re-binned to five forward-split bins}

The binning re-check of 2026-09-05 (technical supplement) found that every released binning
except $\cos\theta_\pi$ is at its resolution limit, and that $\cos\theta_\pi$ admits a fifth
bin in the forward hemisphere. Adopted on 2026-09-06: edges $-1,-0.1,0.35,0.6,0.8,1$
replace $-1,-0.1,0.35,0.75,1$ for the three inclusive extractions (the proton-tagged
$\cos\theta_\pi$ keeps four bins). The four generator predictions were regenerated at the
new edges from the event-level files (\texttt{generator\_predictions/costhpi5}); the
unchanged observables reproduce the previous prediction files to machine precision. The
three universe files were rebuilt (SLURM 3277974), all extractions re-unfolded and the
release regenerated; the D'Agostini cross-check was re-run.

\textbf{Effect.} FHC $\cos\theta_\pi$: $\sigma_\mathrm{int}$ $0.726\to0.674$, closure
$\chi^2$ $0.80/4\to0.68/5$, $A_C$ row sums $0.71$--$0.82\to0.68$--$0.77$,
$s_{\min}/s_1$ $0.27\to0.46$; RHC $0.663\to0.615$, $0.43/4\to0.49/5$; combined
$0.753\to0.658$, $1.05/4\to0.73/5$. The inclusive closure mean moves from $1.018$ to
$1.010$, the range of inclusive integrals to $0.62$--$0.96$, and the D'Agostini-to-Wiener-SVD
ratio to $1.23$--$1.31$ (flat to $\le1.2\%$). No other extraction changes.

\subsection{Documentation-only corrections from the 2026-09-02 review}

None of the following moved a released number; they are recorded because a reader of
Draft~1.1 would otherwise take the printed formulas as the implementation.
(i)~The differential cross-section equation divided by the efficiency in addition to
unfolding with the smearceptance matrix $S=\varepsilon M$; the implementation applies
the correction once, through $S$, and the equation now says so.
(ii)~The response-matrix figure was described as the column-normalised migration
matrix; it is $S$, with columns summing to the efficiency.
(iii)~The muon-candidate cut was described as ``longest qualifying track''; the code
orders the qualifying set by the muon-BDT score.
(iv)~Coherent pion production was described as a background; the signal definition is
topological and includes it.
(v)~The single-exposure normalisation example ($7.63\times10^{20}$~POT) and a
pre-rebuild systematic-breakdown table were replaced by the final three-configuration
values.  (vi)~Flux composition was labelled as integrated above $0.1$~GeV; the
values are for $60$~MeV.  (vii)~TKI angle units in the binning appendix read rad; they
are degrees.  (viii)~A $\theta_\pi$ result was listed as final; none is extracted.
(ix)~The beamline-geometry flux term, quoted at $1.6$--$3.3\%$, is not in the released
systematics configuration and is now stated as absent.

\subsection{External review of Draft~1.5: the two blocking safeguards implemented}

An external reviewer read Draft~1.5 and the supplement together as a package for opening
the inclusive control regions and judged the science mature but the package not yet
approvable, on two explicitly required safeguards that the note itself declared missing,
plus governance and consistency items. Acted on 2026-09-06:
(i)~\textbf{The signal-region-stripped beam-on skim now exists.} \texttt{ProcessNTuples}
with \texttt{XSEC\_CR\_SKIM=1} writes an event only if it fails the signal-region
selection and belongs to at least one control region, suppresses the selection summary,
and stamps the file with a \texttt{XSEC\_CR\_SKIM} marker that \texttt{UniverseMaker}
refuses regardless of \texttt{XSEC\_UNBLIND}; the sideband macros refuse any beam-on
file that has neither the marker nor Monte Carlo truth (\texttt{macros/sb\_guard.h}).
Verified by \texttt{macros/cr\_skim\_check.C} on three inputs, the Run-1 FHC overlay and
the Run-1 and Run-4c beam-off samples, each
processed both ways (recorded in the repository at the tag). No beam-on file was touched.
(ii)~\textbf{The $\pi^0$-definition comparison} with and without the final multiplicity
cut, promised before ratification, was produced from a Monte-Carlo-only reprocess into a
separate tree (\texttt{/data/uboone/processed/sb\_pi0}, SLURM 3277998; new branches
\texttt{sb\_pi0\_final}, \texttt{sb\_nprimtrk}, \texttt{sb\_nnonproton}); results in the
supplement, \S\ref{S-sec:cr_pi0_final}: the cut removes $11$--$14\%$ of the region uniformly across its
components, and the frozen definition is kept.
(iii)~\textbf{Governance.} The run-period rule no longer excludes a period automatically
on a $3\sigma$ excursion: it triggers an investigation, exclusion requires an independent
detector or processing defect and a new frozen release. A fitted $\pi^0$ correction is not
pre-authorised by the $50.1$/$51.0\%$ purity: it may be proposed only after the pre-fit
result is documented and reviewed. What a fitted correction updates in the covariance is
now stated (target-class normalisation only, re-evaluated per universe; flux coherence
retained). The two fitted regions are disjoint by construction, so no joint fit counts an
event twice.
(iv)~\textbf{Consistency.} The supplement's two-bin $p_\pi$ passage still carried a
``$1.13$ in all three'' unfolded-to-truth ratio from a superseded pass beside the current
$0.997/0.981/0.989$; withdrawn, and the $A_C$ diagonals quoted there brought to the
release ($0.53/0.38/0.53$ first-bin diagonals). The note's summary pointed the identity
closure to the supplement's cross-check section instead of its identity appendix; fixed.
``Full covariance'' is now ``complete configured covariance'' with the identified
omissions named beside it; the supplement no longer claims the two documents are
consistent ``by construction''. A boxed scope statement follows the golden-pion
conclusion so it is not read as contradicting the adopted two-bin diagonals; the
supplement's normalisation appendix is retitled as a superseded record; the cosmic region
is introduced as the large-opening-angle, cosmic-enriched region; ``beam-off'', ``EXT''
and ``cosmic data'' are declared synonyms at first use; ``signal contamination'' is
defined as the inclusive signal unless stated otherwise.

\subsection{Second external review (2026-09-07): a stale figure and three captions}

The external reviewer re-read the 2026-09-06 note and supplement. The 1.13 contradiction is
confirmed resolved. Found and corrected on 2026-09-07:
(i)~\textbf{The proton-tagged two-bin $p_\pi$ figure of the supplement was stale}: it
printed integrals of $0.507/0.447/0.490$ against the table's $0.498/0.429/0.462$. The
producing macro (\texttt{macros/ppi2bin\_figs.C}) read the 2026-08-30 sidecars
\texttt{closure\_hists\_xsec1p\_*\_ppi.root} while the tables read the current
\texttt{closure\_hists\_xsec\_ccpi1p\_*\_ppi2bin.root}; the macro now reads the current
sidecars for both families and all nine figures are regenerated (inclusive values
unchanged). The staleness guard could not catch this, since it compares modification
times and the figure postdated the sidecar; it now also reads the integral printed on
each two-bin figure and compares it with the sidecar's.
(ii)~The proton-tagged closure-table caption said the integral ratios ``are the $A_C$
smearing effect''; both numerator and denominator are in $A_C$-smeared space, so the
ratios measure the throw's realisation and any residual bias. Corrected.
(iii)~The bin-integrated two-bin $p_\pi$ table called its uncertainty the data-statistical
term; it is the total from the released covariance ($1.36\times0.030=0.041$). Corrected.
(iv)~The inclusive two-bin figure caption named NEUT ``closest''; the closest generator
changes by bin. Reworded. (v)~The row-sum interpretation is now stated for a flat truth,
with the model dependence acknowledged; ``retains $5\%$ of a bin's normalisation'' became a
row sum of $0.05$; the supplement's ``property of the response, not of the model'' is
qualified to the tested generators. (vi)~The open upper $p_\pi$ bin: the released density
and the conventional width form a matched representation of the bin integral, stated.

\subsection{Stale-detail sweep (2026-09-07)}

Three independent reviewers read the four documents and the data-release index for
statements overtaken by the last week's changes. Nothing numerical moved. Corrected: the
skim verification is now described as three inputs everywhere (Run-1 FHC overlay, Run-1 and
Run-4c beam-off); the remaining ``full covariance'' phrases became ``complete configured
covariance'' (note, talk, release README and index header); the beam-on normalisation
paragraph of the note now states the distinct-exposure rule; the CRT paragraph no longer
says the branches are zero-filled (they are pass-through, populated for Runs 3--5 only);
the supplement's proton-tagged $A_C$ mean row sum for $W_\mathrm{had}$ is $0.61$ (release
value), the pre-fix $\cos\theta_\mu$ generator-suppression passage and the 2026-08-28
combined-detector and spread passages are labelled with their pass, the binning re-check
tables say ``released'' means the 2026-09-04 release, and the $W_\mathrm{had}$ edge scan is
recorded as open (non-gating; note, Limitations); the talk's backup control-region frames
say ``adopted target; ratification requested''; the change log is regrouped (count-level
repairs and review passes as sections), its superseded table gains the cut-and-count
totals, proton-tagged purity, ensemble offsets and FHC beam-off rows, and the note's
count of result-affecting entries is twenty (bundles counted once; the three
presentation- or documentation-only entries excluded).

\subsection{Referee re-review (2026-09-07)}

The external referee re-read the package as a whole. Corrected: the closure-$\chi^2$
discussion said the uncertainty sizing was tested by neither a restricted covariance nor an
ensemble, contradicting the ensemble subsection that follows (reworded); the provenance
sentence intended for the binning re-check tables had landed twice in the total-cross-section
caption (moved); the status matrix now says ``complete configured covariance'' and marks the
proton-tagged differential results as released secondary with binning validation incomplete
($W_\mathrm{had}$ edge scan not run); the $\pi^0$ region's paragraph states the chronology
(cut dropped for a plausible sculpting argument; the comparison found none at the
composition level; the no-cut definition is kept as the larger sample); the $1+B/S$ flux
mechanism is set in a box; the bin-integrated $p_\pi$ representation is declared primary.
Added, on the referee's suggestion: a null calibration of the global control-region test
($2\times10^5$ pseudo-experiments, prediction from the pre-fit covariance, Poisson data):
the nominal $p<0.05$ threshold fails 5.1\% of null experiments and $p<0.01$ 1.1\%, so
the $\chi^2_8$ asymptotics hold for this covariance (finite-universe noise not re-thrown).

\section{Corrections of 2026-09-16/17 (control regions and fit statistics)}

Found while building the joint FHC$+$RHC background fit; all three are documented in full in the
control-region note (its \S5 and \S6), which carries the corrected tables. The released cross
sections are \emph{not} affected by any of them.

\subsection{The sideband Monte Carlo predated the beam-frame correction (numbers moved)}
The sideband-instrumented Monte Carlo \texttt{processed/sb/} was produced 2026-08-11, before the
beam-frame correction of the muon angle (2026-08-19), so it stores $\cos\theta_\mu$ about the
detector $z$ axis while the beam-on control-region skim stores the beam-frame angle. Every
$\cos\theta_\mu$ comparison of the 2026-09-09 control-region note compared two different
conventions. Data over prediction in the CC$0\pi$ region (FHC Runs 4+5) swung from $0.22$ to $2.46$
across the angle bins; with the current Monte Carlo (\texttt{processed/sb\_pi0/}, 2026-09-06) it is
flat at $1.10$--$1.34$. The shape tests move accordingly:

\begin{center}\small
\begin{tabular}{lcc}
\toprule
Normalisation-marginalised shape test & 2026-09-09 & corrected \\
\midrule
CC$0\pi$ $\cos\theta_\mu$, FHC & $350.51/4$ & $3.26/4$ \\
CC$0\pi$ $\cos\theta_\mu$, RHC & $583.58/4$ & $5.48/4$ \\
$\pi^0$ $\cos\theta_\mu$, FHC & $210.36/4$ & $0.59/4$ \\
$\pi^0$ $\cos\theta_\mu$, RHC & $240.85/4$ & $0.49/4$ \\
\bottomrule
\end{tabular}
\end{center}

\noindent The region flags do not depend on the angle, so the yields, the per-run table and the
global normalisation test are unchanged. The $p_\mu$ shape failures ($\pi^0$ in both modes, RHC
CC$0\pi$) survive the correction. The muon-angle shape failure reported in that note is withdrawn.
\texttt{sb\_plots\_data.C}, \texttt{sb\_perrun\_data.C} and \texttt{sb\_protocol\_data.C} now read
\texttt{sb\_pi0/}; the control-region data figures are regenerated (the originals are kept as
\texttt{sideband\_data\_*\_sb20260811.pdf}).

\subsection{Sideband detector variations were not POT-normalised (numbers moved)}
The sideband macros formed the detector-variation shifts as ratios of unscaled event counts. The
knob samples do not have the POT of the detector-variation central-value sample (FHC: WMX $0.913$,
LYdown $0.926$, Recomb2 $1.039$), so that ratio carried the exposure difference as a spurious shift
of about $-8\%$ per knob. The detector term of the control-region covariance:

\begin{center}\small
\begin{tabular}{p{0.40\textwidth}cccc}
\toprule
Detector term (quadrature over 8 knobs) & CC$0\pi$ & $\pi^0$ & cosmic & signal region \\
\midrule
unscaled ratios, FHC / RHC & $20.2$ / $8.4\%$ & $20.1$ / $7.2\%$ & $21.7$ / $30.9\%$ & $21.0$ / $15.4\%$ \\
POT-normalised, event-matched, FHC / RHC & $1.9$ / $2.1\%$ & $0.8$ / $1.8\%$ & $6.5$ / $7.0\%$ & $8.3$ / $8.5\%$ \\
\bottomrule
\end{tabular}
\end{center}

\noindent With the corrected covariance the frozen global normalisation test moves from
$\chi^2 = 8.77/8$ ($p=0.36$, passes) to $69.5/8$ ($p<10^{-3}$, fails); excluding FHC Run~1, whose
exposure is under investigation, it passes again at $5.0/8$ ($p=0.76$). Transfer-factor ratios
(numerator and denominator from the same sample) are unaffected, as is the extraction, which
normalises by \texttt{summed\_pot}. The statistical part of the detector variations is now removed
with an event-matched Poisson bootstrap: the knob samples are re-simulations of the same generated
events ($96.8$--$98.1\%$ matched), so the noise is $3$--$10\times$ smaller than an
independent-sample estimate; in fine bins the median quadrature shift of $8.3\%$ is $7.3\%$ noise.
\texttt{sb\_protocol\_data.C} takes both corrections as switches, on by default; with both off it
reproduces the published tables exactly.

\subsection{Fit statistics: CNP data variance leaves a pull bias (presentation of the new fit only)}
In the background fit, $10^5$-toy ensembles show the combined Neyman--Pearson data variance, which
depends on the observed data, leaves a pull-mean bias of $+0.019$ to $+0.034$ in the
two-dimensional fits (six to ten ensemble standard deviations); it survives Gaussian toys, so it is
not Poisson counting. With the Pearson variance at the nominal prediction the pull mean converges to
zero within $0.005$. The fit default is the Pearson form. No released number is involved.

\subsection{New documentation}
\texttt{report/newobs\_note.tex} records the $0.50$ migration-diagonal study (no rebinning adopted:
all twenty-six candidates are rejected by the $A_C$ conditioning), the new proton-tagged observables
$\theta_p$ and $\theta_{\pi p}$ with their reprocess and its bit-for-bit gate, and the
two-dimensional programme. The control-region note gains its \S5 (corrections) and \S6 (joint
FHC$+$RHC background fit).

\section{Corrections and additions of 2026-09-18}

\subsection{\texorpdfstring{$\theta_p$}{theta p} adopts RHC-derived edges}
\begin{sloppypar}
The maximin binning scan had only ever been run on the FHC chain, because
\texttt{macros/tki\_binning\_scan.C} chained the FHC files unconditionally; it now takes a beam
argument, FHC by default so every published screen reproduces. Run on RHC it returns different edges,
and at the \emph{same} four bins those edges condition three to four times better in both usable
configurations (FHC $s_{\min}/s_1$ $0.082\to0.383$, combined $0.125\to0.408$). $\theta_p$ therefore
adopts $0,\,0.659,\,1.068,\,1.539,\,\pi$~rad. Coarsening to three bins does \emph{not} rescue the
rank-deficient RHC response ($0.002\to0.050$), so $\theta_p$ is quoted for FHC and the combined
configuration only. The price of the change is stated with it: the data-statistical term rises from
$12.3\%$ to $16.6\%$ (FHC) and $9.7\%$ to $14.7\%$ (combined), because better conditioning means
less of the truth smeared into the measurement, not a smaller error bar. Configurations, generator
predictions and all three extractions were rebuilt; regenerating the predictions left $240$
histograms byte-identical and changed only \texttt{thetap\_fte}. Superseded outputs are kept as
\texttt{*\_thetap\_fhcedges}. No released observable is affected.
\end{sloppypar}

\subsection{The proton candidate, and where its cost actually appears}
On selected signal events the reconstructed proton candidate is the leading true proton $81.7\%$ of
the time, some true proton $87.4\%$, and not a proton at all $12.6\%$ (the $59\%$ recorded during
the screen predates the proton-tagged reprocess and is withdrawn). A wrong pick randomises the angle
rather than smearing it, so it widens the core resolution by only a quarter to a third while
supplying $74\%$ of the residuals beyond $0.3$~rad in $\theta_p$ and $51\%$ in $\theta_{\pi p}$:
candidate purity dominates bin-to-bin \emph{migration}, which is the mechanism behind $\theta_p$'s
marginal diagonal and its rank-deficient RHC response. No released number changes.

\subsection{Far sidebands, and a ratio that does not move}
Adding one cut per control region drives the CC$0\pi$ signal fraction from $8.6\%$ to $2.2\%$ while
keeping $58\%$ of the events and raising the CC$0\pi$ content of the background from $71\%$ to
$85\%$. The data-to-prediction ratio is flat under the far cuts ($1.12\to1.13$; $\pi^0$
$1.20\to1.23$), so the common excess of data over prediction is not signal contamination.
\emph{Corrected the same evening:} the first version of this study reported an $8$--$16\%$
\emph{deficit} because \texttt{bkgfit/farsb.py} summed the prediction over the FHC and RHC Run-2
periods, which have simulation but no beam-on file ($\sim20\%$ of the prediction), and because the
FHC Run-1 exposure was still wrong; with both fixed the ratios are $1.1$--$1.2$, and the cosmic region
falls to unity as its beam-off share rises, locating the excess in the neutrino simulation.

\subsection{Corrections to claims in the documents}
\begin{itemize}\itemsep2pt
  \item The technical supplement stated that the released proton-tagged binnings ``were not
    re-tested''. They have been, at both $0.68$ and a relaxed $0.50$ diagonal: twenty-six candidates
    pass the relaxed diagonal and all fail on $A_C$ conditioning. Only a two-bin $W_\mathrm{had}$ is
    healthy, and it was not adopted.
  \item The analysis note carried the request to open the control regions with no record that they
    were opened on 2026-09-09. It now records the outcome, including that the sign-off checklist
    was not revised at opening.
  \item \sloppy \texttt{norm\_manifest.py} reported a spurious \textsc{unfold prescription drift} and
    rewrote the tracked manifest on every run: the one-bin total extraction is deliberately
    unregularised, so the two classes are now audited separately.
\end{itemize}

\section{Corrections of 2026-09-19}

\subsection{The FHC Run-1 exposure: confirmed and adopted}
The FHC Run-1 control-region deficit (fitted exposure factor $0.520\pm0.020$; every region at
$0.60$--$0.72$ against unity elsewhere) had a specific cause: the beam-on sample was recorded
\emph{without} the open trigger, $2.192\times10^{20}$~POT with $5\,748\,692$ triggers, while every
configuration carried the open-trigger pair, $3.283\times10^{20}$ with $9\,846\,635$. Six pairings
were tested against the four regions, including a third trigger count found in the repository; only
the non-open pair brings all four to unity ($0.970$, $1.090$, $1.380$ on $40$ events, $0.948$), and
the beam-off-dominated cosmic region discriminates the gate count on its own. The pairing was then
confirmed against the beam bookkeeping.

Adopted in every configuration (17 file-properties lists, the throw macros, the gate counts in the
analysis code, the normalisation audit); FHC Run~1 alone re-thrown; FHC exposure $8.857\to
7.766\times10^{20}$, combined $1.9939\to1.8848\times10^{21}$. The FHC and combined extractions are
rebuilt on it (universe files, unfolding, covariances, release, figures and tables). On the blind
fake data the exposure cancels, because the pseudo-data are thrown at the assumed exposure and
numerator and denominator move together, and the rebuilt integrals differ from the previous release only
through the new FHC Run-1 Poisson realisation and the changed run weighting (median ratio $0.99$;
individual bins within their statistical uncertainties). On beam-on data, where the event count is
fixed, the same correction raises $\sigma_{\rm FHC}$ by $14.0\%$ and $\sigma_{\rm comb}$ by
$5.8\%$ relative to the old bookkeeping; RHC is unaffected. The combined flux constant is
re-weighted to the corrected POT ($6.60865\to6.59691\times10^{-10}$~cm$^{-2}$/POT). \texttt{bkgfit/run1\_trigger.py} is the record of the
evidence.

\subsection{The rebuild on the corrected exposure landed; what moved and what did not}
All 34 FHC and combined universe files were rebuilt on the corrected exposure (SLURM 3351878),
promoted (the previous files are kept in \texttt{univmake\_backup\_job3351878}), and the full chain
re-run: 51 re-unfolds, covariance harvest, release export, figures, counting and every
release-derived table. The one-bin totals, which the rebuild manifest had omitted, were rebuilt
separately (SLURM 3352400). Because the pseudo-data are thrown at the assumed exposure, the
exposure cancels in a fake-data extraction and nothing moved by $14\%$: the rebuilt integrals are
$0.95$--$1.05$ of the previous release for the inclusive observables (median $0.99$) and within
$\pm30\%$ in the sparsest proton-tagged bins, all inside their statistical uncertainties, from the
new FHC Run-1 Poisson realisation plus the changed run weighting (Run~1 is now $28\%$ of FHC,
not $37\%$). The closure mean is unchanged ($1.011$ inclusive, $1.012$ proton-tagged). The
$+14\%$ / $+5.8\%$ applies to a beam-on extraction at fixed event count, and the documents now say
so.

\subsection{Three further places the correction had not reached}
\begin{itemize}\itemsep2pt
  \item \textbf{The combined flux constant} was still the $8.857:11.082$ POT-weighted mean of the
    FHC and RHC constants ($6.60865\times10^{-10}$); re-weighted to $7.766:11.082$ it is
    $6.59691\times10^{-10}$~cm$^{-2}$/POT ($-0.18\%$), applied to all 50 combined configurations
    and audited by \texttt{norm\_manifest.py}; the combined extractions were re-unfolded.
  \item \textbf{Hard-coded per-run scales.} Some twenty standalone macros and tools (cut-flow
    yields, the counting cross-check, binning screens, every sideband macro, the background-fit
    templates, \texttt{farsb.py}) carry the per-run data/MC POT scale, the full-exposure dirt scale,
    the beam-on trigger counts or the combined flux as literals. FHC Run~1 was $0.14101$
    ($3.283/23.282$) in all of them; it is $0.09415$ now, with the dirt scale $0.092402\to0.081020$
    where it stands for the full FHC exposure and the pooled-EXT factors of the counting macro
    recomputed. The framework itself takes all of this from the file lists and was never affected.
    The cut-flow tables and the counting log are regenerated from the corrected macros; the
    background-fit templates and the sideband protocol are being re-run on them.
  \item \textbf{Exposure-derived numbers in the documents} (POT tables, integrated flux and
    conversion factors, the combined $\nu_\mu/\bar\nu_\mu$ flux row, the cut-flow headers).
\end{itemize}

\subsection{Background fit and frozen protocol re-run on the corrected exposure}
The background-fit templates (\texttt{bkgfit\_templates.C}) and the frozen control-region protocol
(\texttt{sb\_protocol\_data.C}) carried the wrong Run-1 scale; both are re-run. The frozen global test
now \textbf{passes}, $\chi^2=8.1/8$ ($p=0.42$; it was $69.5/8$), and the all-period totals fit needs no
parameter ($5.7/6$; CC$0\pi$ free $0.949\pm0.190$; the ``no FHC Run~1'' results are unchanged). A
Run-1-specific residual survives: every other period sits $12$--$25\%$ above the pre-fit prediction,
FHC Run~1 at unity, so a free Run-1 factor still fits $0.869\pm0.022$ ($\Delta\chi^2=37$; a free
factor on any other period gains at most $15$). The exposure correction explains $0.52\to0.87$, not
$0.52\to1$. Its cause is not identified and it is a new open item for unblinding; the CR note carries
the per-period table. The sign of the ``residual deficit'' statements has changed with the
correction: on the corrected exposure the regions sit \emph{above} the prediction, and the
far-sideband study is being re-run.

\subsection{The MCS momentum-scale bound is withdrawn and redone on the data side}
The $\pm5\%$/$\pm10\%$ bound quoted on 2026-09-19 scaled the MCS momentum through the
\emph{bin configuration}, which applies to the simulation and to the fake data alike. A common
scale cancels in the extraction, since response and data move together, so the $0.07\sigma$ and
$0.19\sigma$ measured only the regulariser's response to a remapped variable, not the systematic.
A momentum-scale error is a data-versus-simulation mismatch, so the variation is now applied to
the fake data alone (\texttt{macros/mcs\_scale\_fakedata.C}: every MCS-measured muon, $63\%$ of the
selected sample, has its momentum scaled by $1\pm0.05$; the simulation, the response and the
selection flag are untouched) and $d\sigma/dp_\mu$ is re-extracted with the nominal response
(SLURM 3352423). Result: the half-difference is $8$--$13\%$ of the bin value in the four lowest bins,
$2\%$ in the fifth and $8$--$10\%$ with the opposite sign in the two highest, i.e.\ up to $0.34\sigma$
of the quoted per-bin uncertainty (the withdrawn common-scale test had given $0.07\sigma$); the integral
moves $-6.8\%$ / $+2.1\%$. As recommended by the review it is released as a fully correlated term
for FHC $p_\mu$ (\texttt{cov\_MCSscale.txt}, \texttt{cov\_total\_plusMCS.txt}; per-bin totals
$38$--$58\%\to39$--$58\%$); propagation inside the framework and to the other configurations and
observables is open.

\subsection{Proton-tagged alternative-model closure}
The reweighted-model fake data (GENIE multisim universe 545; $\Delta\to N\pi$ angular unisim) are
thrown for the proton-tagged sample as well (\texttt{macros/throw\_altmodel\_w.C}, same universes
and seeds) and $W_\mathrm{had}$, the four two-bin TKI variables and $\cos\theta_\pi$ are
re-extracted with the nominal response (SLURM 3352430), one of the six prerequisites listed for
the proton-tagged products. Outcome: nothing beyond $0.9\sigma$ under the $\Delta\to N\pi$
variation; under the coherent GENIE variation the low-imbalance $\delta p_T$ and $p_n$ bins come
back $35\%$ below the alternative truth ($-1.14\sigma$, $-1.16\sigma$) while that truth moved
$1\%$, i.e.\ the shift enters through the response and background under the varied model. Larger
than the inclusive maximum ($0.93\sigma$); recorded in the proton-tagged note as a model
sensitivity of those bins at the level of their uncertainty.

\subsection{Study extractions rebuilt on the corrected exposure (2026-09-20)}
The proton-angle products ($\theta_p$, $\theta_{\pi p}$; FHC and combined, live tree), the sixteen
FHC and combined two-dimensional pairs and the surviving $\theta_p$ binning candidates were built
before the Run-1 correction and are rebuilt on it (SLURM 3352599, 24 extractions; the previous
live $\theta_p$/$\theta_{\pi p}$ files are kept in \texttt{univmake\_backup\_study\_pre\_run1fix}).
As for the release, the conditioning is a property of the response and barely moves; the
throw-dependent numbers do: $\theta_p\times\delta p_T$ conditions at $0.611$ combined (was $0.617$),
$\cos\theta_\pi\times\cos\theta_\mu$ at $0.564$ ($0.563$), the two-dimensional chain reproduces the
one-bin total to a median of $1.5\%$ ($0.7\%$), $\theta_p$ conditions at $0.325$/$0.409$ FHC/combined
($0.383$/$0.408$), and the three-bin $\theta_p$ alternative at $0.12$/$0.05$, combined $0.49$. The
2D figures, the proton-angle figure, the inventory and information-content tables and the
proton-tagged note's 2D table are regenerated; the conclusions (two pairs usable, combined only;
$\theta_p$ FHC and combined only) are unchanged. The inclusive FHC alternative-model fake data,
thrown on 2026-09-18 at the old exposure, are re-thrown at the corrected one and the eight FHC
extractions re-run (SLURM 3352603): the largest bias is now $0.93\sigma$ (a sparse $\cos\theta_\mu$
bin under the angular variation; was $0.81\sigma$), the near-threshold $p_\pi$ region returns $0.40$
of the alternative truth (was $0.415$) and the sparse top $p_\mu$ bin $0.58$ (was $1.475$: the
tail bin is dominated by its own realisation); the populated bins are robust as before.

\section{Corrections of 2026-09-24}

\subsection{Ensemble coverage: ten members rebuilt from truncated fake data (numbers moved)}
Completing the six ensembles to $100$ members raised the RHC-total statistical pull width from
$1.10$ ($46$ members) to $1.28$, with the ensemble mean $-1.5\%$ ($2.9\sigma$) low. The cause was
a race condition in the ensemble machinery, not the extraction: the throw file of a member is
named by beam mode and throw number only (\texttt{ens/fakedata\_rhc\_run2\_t57.root}), and every
ensemble job re-throws it. The combined and RHC-total ensembles started the same throw numbers
at the same moment (SLURM 3352020/3352022, 2026-09-20 17:10), so one job read the file while the
other was recreating it and silently saw part of the tree. A check of every member (events read by
univmake against entries on disk, $2800$ inputs) found RHC-total members $56$, $57$, $59$, $60$ with
Run~2 read at $47$--$72\%$ and FHC $\cos\theta_\mu$ member $26$ at $43\%$; five further RHC-total
inputs differ by one entry, which proved to be an out-of-range truth event in the one-bin
binning and not a truncation. The throw macros now write to a temporary file and rename it on
completion; the affected members were quarantined (\texttt{ens/corrupt\_race\_2026-09-20/}) and
re-run (SLURM 3427769/3427770). Corrected results, all at $100$ members: statistical pull widths
$0.92$--$1.07$ (RHC total $1.28\to0.99$), $68\%$ coverage $64$--$74\%$, ensemble means $0.4$--$1.1\%$
low at $0.5$--$2.0\sigma$ on the mean (the $2.0\sigma$ is the RHC total, $-0.81\%$). The coverage
tables of the note and the supplement and \texttt{data\_release/ensemble\_2026-09-20.tsv} are
regenerated.

\subsection{MCS momentum-scale term extended to RHC and combined $p_\mu$}
The data-side variation of 2026-09-19 is now run for the other two configurations
(\texttt{macros/mcs\_scale\_fakedata.C} takes the beam mode; RHC fake data $\pm5\%$ on the $60\%$ of
selected muons measured by MCS; the combined configuration scales FHC and RHC together; SLURM
3427764). Combined: half-differences $5$--$10\%$ in the five lowest bins, $2\%$ and $14\%$ in the two
highest, at most $0.35\sigma$; integral $-2.5\%$/$+2.4\%$. RHC: $6$--$15\%$ ($\le0.32\sigma$) in
the six lowest bins but $+21\%$/$-44\%$ ($0.85\sigma$) in the open top bin; integral
$-0.8\%$/$-6.8\%$. Released as \texttt{cov\_MCSscale.txt} and \texttt{cov\_total\_plusMCS.txt}
under \texttt{incl\_RHCFULL\_pmu} and \texttt{incl\_COMB\_pmu} (\texttt{report/tools/mcs\_eval.py CFG}).
The RHC top bin is the one place where the term is not small against the quoted uncertainty.
Propagation inside the framework and to the other observables remains open.

\subsection{MCS momentum-scale term is zero for the other inclusive observables}
The same data-side $\pm5\%$ variation was extracted for $\cos\theta_\mu$, $\cos\theta_\pi$,
$p_\pi$ (two regions) and the one-bin total (SLURM 3427873, stopped after 9 of 36 jobs). Every result
is identical to the nominal one (ratio $1.000000$ in every bin): the variation changes only the
reconstructed muon momentum, which no selection cut and no other inclusive observable uses. The term
is therefore exactly zero for all inclusive observables except $p_\mu$. It remains to be evaluated
for the proton-tagged transverse-kinematic observables, which use the muon momentum vector.

\subsection{Values corrected while rewriting the analysis note (numbers moved)}
The analysis note was rewritten in the order of the analysis, without release history, and the
following values that the exposure correction of 2026-09-19 had not reached were corrected:
\begin{itemize}\itemsep1pt
  \item Combined generator predictions: \texttt{combine\_comb\_fte.C} and \texttt{combine\_ext.C}
    weighted the modes with the old FHC exposure ($8.857\times10^{20}$), FHC fraction $0.458$ instead
    of $0.425$. The combined generator files were regenerated (their predictions move by $-0.6\%$
    inclusive and $-0.8\%$ proton-tagged; the old files are in
    \texttt{generator\_predictions/comb\_backup\_pot8857/}) and the result set re-unfolded; only the
    combined generator curves and $\chi^2$ change.
  \item \texttt{data\_release/ext\_gates.tsv} and the beam-off figure still carried FHC Run~1 at
    $3.283\times10^{20}$~POT and $9\,846\,635$ gates (FHC beam-off $30.2$ events); both are now
    generated by \texttt{report/tools/ext\_gates\_table.py} from the live file lists (FHC $23.2$, RHC
    $23.2$, combined $46.4$ events).
  \item Note: the combined data POT of the exposure table ($19.939\to18.848\times10^{20}$), the
    combined flux composition ($53.0/44.4\to52.4/45.1\%$), the background composition at the final cut
    (old exposure), the control-region yields and pseudo-data ratios for FHC and combined (regenerated
    with \texttt{sideband\_compare.C}), the two-region $p_\pi$ table (FHC $0.059/0.764\to0.037/0.827$;
    now generated by \texttt{rebuild\_tables.py}), the $A_C$ row-sum table (now generated; FHC
    $\cos\theta_\mu$ $0.17$--$0.39\to0.22$--$0.42$), the flux-term arithmetic ($B/S=560/833$), the
    one-bin breakdown ($38.5\oplus4.3\to38.0\oplus4.5\%$), the combined mode weights ($0.80{:}1\to0.70{:}1$)
    and the FHC beam-on control-region figure (regenerated on the corrected exposure).
  \item \texttt{tables/infocontent.tex} read hand-copied uncertainties of an older release
    ($p_\pi$ $53.9\%$ against $73.0\%$); it now reads the systematics dumps.
  \item The glossary defined the migration diagonal as a row-normalised fraction; it is the diagonal
    of the column-normalised migration matrix, as used throughout.
\end{itemize}

\section{Additions of 2026-09-25}

\subsection{Statistical coverage of the proton-tagged family}
Three $100$-member fake-data ensembles now cover the proton-tagged family
(SLURM 3427889--91, job script \texttt{slurm\_ensemble\_1p.sbatch}): $W_{\rm had}$ and the two-bin
$\delta p_T$ in FHC, and the two-bin $p_n$ in the combined configuration. The throws are written to
a temporary file and renamed, so the race of 2026-09-20 cannot recur, and every member was checked
(events read against entries on disk, $1600$ inputs, all equal). Statistical pull widths
$1.02$, $1.08$, $0.99$; $68\%$ coverage $68.7$, $62.0$, $64.5\%$; ensemble-mean integrals
$+0.02$, $-0.65$, $-0.32\%$ from the reference ($\le0.8\sigma$ on the mean). The sparse first
$W_{\rm had}$ bin is the least regular (mean $-0.21$, width $1.14$), within two standard errors.
Prerequisite~3 of the proton-tagged note is met for these three observables
(its section ``Statistical coverage''). The table is generated by
\texttt{ensemble\_tables\_1p.py} and the numbers are released in
\texttt{ensemble\_1p\_2026-09-25.tsv}.

\section{Binning review of 2026-09-26}

\subsection{The 0.68 and 0.50 criteria applied uniformly (study; no binning changed)}
Every observable of both families was re-binned on the corrected exposure under the nominal $0.68$
and a relaxed $0.50$ column-normalised migration criterion, applied in FHC and RHC, with an event
floor ($70$ selected FHC signal events per inclusive bin, $50$ per proton-tagged bin). At every bin
count up to twelve the partition with the largest worst-bin diagonal was found by dynamic
programming (\texttt{report/tools/maximin\_binning.py}), and $141$ candidates were extracted in full
in FHC, RHC and combined (SLURM 3428213, 3428222, 3428223, 3433027, 3433033). The shape each binning
retains is measured by the Wiener filter factors, the eigenvalues of $A_C$, and their sum, the degrees
of freedom for signal; a binning is selected as the fewest bins reaching $90\%$ of the largest
admissible sum (\texttt{report/tools/binreview\_tables.py}, \texttt{data\_release/binreview\_choices\_2026-09-26.tsv}).
The supplement section (\S\ref{S-sec:binreview}) gives the tables, the figure and the observable-by-observable
findings. The released $p_\mu$ (worst diagonal $0.33$) and $W_{\rm had}$ ($0.09$) fail both criteria;
$\cos\theta_\mu$ ($0.67$) and $\theta_\mu$ ($0.65$) fail $0.68$ and pass it with optimised edges at the
same five bins; no two-bin $\delta\alpha_T$ passes $0.68$; the $\theta_p$ and $\theta_{\pi p}$ study
binnings hold $21$ and $42$ events in their sparsest bin. Selected under $0.68$ and $0.50$: $p_\mu$ four
and six bins, $p_\pi$ two and three, $\cos\theta_\mu$ seven and eight, $\theta_\mu$ eight and ten,
$\cos\theta_\pi$ seven, $\theta_{\mu\pi}$ five to ten (equivalent), $W_{\rm had}$ two, $W_{\pi p}$ one and
two, $\delta p_T$ and $p_n$ two, $\delta\alpha_T$ one and two, $\delta\phi_T$ two and three, $\theta_p$
four, $\theta_{\pi p}$ three and four. The same partition retains up to one degree of freedom more as
$\cos\theta_\mu$ than as $\theta_\mu$, because the regulariser acts on the density. The review
recommended $0.68$ for every observable; the decision, taken the same day, is recorded in the next
section. A single-throw closure cannot establish the coverage of a new binning, so each adopted one
needs its own ensemble.

\subsection{Statements corrected by the review (text moved)}
Four statements about the released binnings were wrong and are corrected. The note's appendix on
alternative binnings said that every binning is the coarsest that meets the criterion and that the
conditioning of $A_C$, not the diagonal, limits the binning: the released $p_\mu$ and
$W_{\rm had}$ fail both criteria, $\cos\theta_\mu$ and $\theta_\mu$ fail $0.68$, and the threshold of
about $0.1$ on $s_{\min}/s_1$ used to reject candidates is not met by most released binnings. The
note's appendix on distributions not reported said that the $W_{\rm had}$ response is adequate at
six bins: all six fail $0.68$ and five fail $0.50$. The supplement's re-check said that
$\cos\theta_\mu$ and $\theta_\mu$ cannot go finer: with optimised edges both admit more bins, and
their released five bins fail $0.68$ where optimised five-bin edges pass. The RHC-derived $\theta_p$
edges were presented as an improvement: they leave $21$ selected events in the last bin. The
proton-tagged note's binning prerequisite now states the $W_{\rm had}$ finding, and the note's
unblinding requirements include the choice of criterion and binnings.

\section{Binnings of the 0.50 criterion adopted (2026-09-26)}

\subsection{The decision and the binnings}
The $0.50$ migration criterion was adopted for both families, with the binnings selected by the
rule of the review (supplement, \S\ref{S-sec:binreview}) and four exceptions: $\cos\theta_\pi$ keeps
its five bins, $\theta_{\mu\pi}$ its five, $\delta\alpha_T$ its split at $120^\circ$, and the
proton-tagged $p_\pi$ its two regions, since three fail $0.50$ in that sample. The binnings are
$p_\mu$ six bins ($0.15$, $0.30$, $0.425$, $0.575$, $0.825$, $1.25$, open) and $\cos\theta_\mu$ eight
($-1$, $0.4$, $0.625$, $0.75$, $0.825$, $0.9$, $0.95$, $0.975$, $1$) in both families; the inclusive
$p_\pi$ in three regions ($0.175$, $0.205$, $0.26$, open); $W_{\rm had}$ two regions split at
$1.18$~GeV/$c^2$; $\delta\phi_T$ three bins ($0$, $22.5$, $75$, $180^\circ$); $\theta_p$ four bins
($0$, $0.44$, $0.785$, $1.162$, $\pi$) and $\theta_{\pi p}$ four ($0$, $0.911$, $1.414$, $1.916$, $\pi$)
at their maximin edges. $d\sigma/d\theta_\mu$ is no longer reported: a partition has the same
diagonals in $\theta_\mu$ and $\cos\theta_\mu$, and the cosine retains up to one degree of freedom
more. The exemption of $p_\mu$ from the criterion ends.

\subsection{Proton momentum added to the proton-tagged family}
The momentum of the leading proton, $p_p$, is a new proton-tagged observable. Its candidates, two to
five bins between the signal threshold of $0.3$~GeV/$c$ and an open top, were extracted in the three
configurations (SLURM 3433088); the rule selects three bins split at $0.42$ and $0.60$~GeV/$c$ (worst
diagonal $0.74/0.73$, trace $1.49/1.40$), and two bins under $0.68$. Its generator predictions were
built from the event-level files. The release now holds fifty-one extractions: fifteen inclusive and
thirty-six proton-tagged.

\subsection{Chain}
{\sloppy Generator predictions at the new edges were built from the event-level files
(\texttt{generator\_predictions/gen2d}; the histograms of the previous binnings are unchanged and the
integrals of every new axis equal those of the old ones). The $27$ universe files of the new
binnings (SLURM 3433061) and the $p_p$ candidates were installed by \texttt{promote\_050.sh}, which
backed up the replaced configurations and universe files
(\texttt{configs/backup\_pre050\_20260926}, \texttt{univmake\_backup\_pre050\_20260926}), archived the
products of the retired extractions (two-region inclusive $p_\pi$, $\theta_\mu$, two-bin
$\delta\phi_T$) and regenerated the release; the retired figures moved to
\texttt{figures/retired\_pre050\_20260926}. The MCS term was re-evaluated for the six $p_\mu$ bins
(largest $0.56\sigma$, in the open top bin in FHC). The alternative-model closure was repeated for
every rebinned observable in FHC (SLURM 3433100, 3433125) and for $p_\mu$, $\cos\theta_\mu$ and
$p_\pi$ in RHC (SLURM 3433342): the largest inclusive bias rose from $0.93\sigma$ to $1.18\sigma$
(the third FHC $\cos\theta_\mu$ bin under the angular variation, a bin with $114\%$ uncertainty),
RHC stays below $0.83\sigma$, and $p_p$ below $0.37\sigma$. Also repeated: the D'Agostini cross-check on the five inclusive observables
(integrals flat in the observable to $1\%$ and equal to the one-bin totals), and the five ensembles
of rebinned extractions (FHC $p_\mu$, FHC and RHC $\cos\theta_\mu$, combined $p_\pi$, proton-tagged
FHC $W_{\rm had}$; SLURM 3433218, 3433219, 3433290, 3433291, 3433292; $500$ members, none failed,
no truncated read): statistical pull widths $1.07$, $0.97$, $1.05$, $0.95$ and $0.94$, integral
offsets $-0.4$ to $-1.1\%$ at $0.4$--$1.8\sigma$ of the mean
(\texttt{data\_release/ensemble\_2026-09-26.tsv}, \texttt{ensemble\_1p\_2026-09-26.tsv}).\par}

\subsection{Defects found during the adoption}
\begin{itemize}\itemsep2pt
  \item \textbf{Ensemble throws on NFS.} Writing each throw to a temporary name and renaming it
    over the final one (2026-09-24) did not stop the race: on NFS the rename unlinks the old file
    under a reader on another node, which then fails with a short read (two FHC $p_\mu$ members
    within two minutes of five arrays starting together). The arrays were cancelled after six
    minutes. A member's throw is now created once by \texttt{link(2)} from a closed temporary,
    stamped with its seed and POT scale, reused when the stamp matches and never replaced while it
    may be open (\texttt{macros/throw\_guard.h}, \texttt{ens/throws/}). The FHC and RHC arrays of one
    observable also shared a work directory, a universe log and a file list per member; these are
    now per configuration.
  \item \textbf{Response figure of $p_\pi$.} The figure script skipped the $p_\pi$ response, so the
    note showed the two-region response of an earlier release; it now writes
    \texttt{fw\_resp\_ppi3bin} and \texttt{fw\_eff\_ppi3bin}.
  \item \textbf{MCS tool.} It read the MCS extractions of the retired binnings for
    $\cos\theta_\mu$; it now checks whether the reconstructed bins read the muon momentum and writes
    the identically zero term otherwise.
  \item \textbf{Proton kinetic energy.} The proton candidate's kinetic energy is the range-based
    value, an exact function of the track length, not a calorimetric one; the notes said otherwise.
  \item \textbf{Reproduction precision.} The external reproduction of the released curves was quoted
    at $1.2\times10^{-10}$; from the released text files it is a relative $4\times10^{-8}$, the
    precision of those files.
  \item \textbf{Hand-maintained tables.} The note's table of bin edges and the supplement's tables of
    $A_C$ row sums and proton-tagged binnings carried values of earlier releases; the first two are
    now generated (\texttt{binning\_table.py}, \texttt{ac\_rowsums\_table.py}).
  \item \textbf{Combined-truth test.} The check of the combined truth against the average of its
    inputs used POT weights and $A_C$-smeared integrals; with fluence weights on the one-bin truths
    the ratio is $1.000$ (inclusive) and $1.001$ (proton-tagged).
  \item \textbf{Top bin closed at a physical boundary.} The last true bin of $\cos\theta_\mu$ and
    $\cos\theta_\pi$ is written as ``$<1.000$'' (likewise $\delta\phi_T$ at $180^\circ$ and
    $\theta_{\pi p}$ at $3.142$~rad), so an event exactly at the boundary falls in no true bin. Five of
    the $145\,946$ signal events of the fourteen overlays have $\cos\theta_\mu=1.0$ exactly (two of
    them selected), none has $\cos\theta_\pi=1.0$, and no ensemble member is short by more than this.
    The effect is $0.003\%$ of the signal and is left in the release; the configuration writer now
    leaves such a bin open at the top for the next rebuild, and the ensemble entry-count scan
    (\texttt{report/tools/ensemble\_entry\_scan.py}) reports a shortfall of one or two entries as this
    case rather than as a truncated read. The one-entry-short members flagged on 2026-09-24 were this
    case, not corrupt reads.
\end{itemize}

\section{Response to the collaboration review (2026-09-27)}

\subsection{Decisions}
\begin{itemize}\itemsep1pt
  \item Four status labels only: \emph{Primary result}, \emph{Secondary result proposed for approval},
    \emph{Validation-only product}, \emph{Not reported}. The status of every product is kept in one
    file, \texttt{report/status.tsv}, from which the status tables of both notes and the status column
    of \texttt{data\_release/index\_extractions.tsv} are generated.
  \item Proposed for approval as secondary results: the proton-tagged one-bin total; the proton-tagged
    $\theta_p$ and $\theta_{\pi p}$ (FHC, RHC, combined); the two-dimensional $\theta_p\times\delta p_T$
    (proton-tagged sample) and $\cos\theta_\pi\times\cos\theta_\mu$ (inclusive sample), combined
    configuration only. Earlier drafts called $\theta_p$, $\theta_{\pi p}$ and the two pairs ``proposed
    secondary'' in some places and validation-only in others. The inclusive pair moved from the
    proton-tagged note to the analysis note.
  \item The MCS muon-momentum-scale term is part of the official covariance of every result
    (\texttt{cov\_total\_plusMCS.txt}); every uncertainty and model-comparison table is regenerated with
    it (\texttt{report/tools/official\_cov.py}). FHC $p_\mu$ moves from $39.3\%$ to $40.6\%$ systematic
    and from $41.2\%$ to $42.4\%$ total; the model $\chi^2$ of $p_\mu$ decrease (FHC GENIE $2.73\to2.24$).
    The closure $\chi^2$ and the robustness biases stay with the framework covariance.
  \item The independent-generator closure is defined: a validated NuWro overlay, every bin holding at
    least $10\%$ of the largest truth within $1\sigma$, every integral within its uncertainty (U3).
\end{itemize}

\subsection{Release}
\begin{itemize}\itemsep1pt
  \item $\theta_p$ and $\theta_{\pi p}$ (three configurations) and the two combined two-dimensional
    results were added to the data release with their covariances, $A_C$ and curves
    (\texttt{xsec\_analyzer/slurm/harvest\_proposed\_local.sh}, \texttt{report/tools/export\_2d.py}). The
    re-harvest reproduced the unfolded results, truths and $A_C$ exactly; the combined generator curves
    of the two-dimensional results moved by $0.3$--$1.1\%$, because their files predated the correction
    of the combined generator weights ($8.857\to7.766\times10^{20}$ FHC exposure).
  \item The MCS term of the new entries is exactly zero for $\theta_p$, $\theta_{\pi p}$ and
    $\cos\theta_\pi\times\cos\theta_\mu$; for $\theta_p\times\delta p_T$, which uses the muon momentum
    through $\delta p_T$, it was extracted (SLURM 3433847): at most $16\%$ of a cell and $0.20\sigma$,
    integral unchanged.
  \item The release README states the $60$~MeV flux convention, the status labels, the official
    covariance and the file schema.
\end{itemize}

\subsection{Documents}
\begin{itemize}\itemsep1pt
  \item Both notes open with the same decision box and a status table; the inclusive note has a
    measurement-definition table, a section on the migration criterion and binning rule, the validation
    ladder, and a numbered prerequisite table with criteria, owners, evidence and blocking levels, shared
    verbatim by the three documents (\texttt{report/shared/}).
  \item The proton-tagged note is restructured: decision and status, scientific purpose, selection,
    observables, extraction, validation, proposed results, validation-only products, products not
    reported, prerequisites. Reconstruction validity (P4) and model sensitivity (P3) are separate
    prerequisites.
  \item The figures use physics notation, the open-ended last bins are hatched and labelled, and the
    six-panel proton-tagged figures are split in two.
  \item The proton-tagged W$_\mathrm{had}$ and TKI observables had no generator predictions in RHC and
    combined, and the muon and pion observables of the proton-tagged sample none in any configuration.
    The RHC and combined W$_\mathrm{had}$ and TKI predictions existed but were not referenced by the
    configurations; the muon and pion ones were produced with new readers that reuse the proton-tagged
    signal definition (\texttt{generator\_predictions/gen2d/*\_1p\_inc*}, \texttt{regen\_1pinc.sh}).
    The $25$ extractions were unfolded again: unfolded results, truths, $A_C$ and every covariance are
    bit-identical, and only the generator curves were added.
  \item Terminology: pseudo-data, control region, combined configuration
    (\nolinkurl{report/tools/terminology.py}). A cross-document check,
    \nolinkurl{report/tools/check_documents.py}, enforces the shared sources, the status vocabulary and the
    adopted exposures, fluxes and target count.
\end{itemize}

\section{Muon momentum-scale term for the proton-tagged observables (2026-09-27)}

The data-side MCS momentum scale ($\pm5\%$) is evaluated for the proton-tagged observables that use
the muon momentum: $p_\mu$, $W_\mathrm{had}$, $\delta p_T$, $\delta\alpha_T$ and $p_n$, in FHC, RHC and
combined (SLURM 3433817, \texttt{xsec\_analyzer/run\_mcs\_1p.sh}). The scaled fake data
(\texttt{macros/mcs\_scale\_fakedata\_1p.C}) recompute $W_\mathrm{had}$ and the TKI variables from the
scaled muon; with a scale of one they reproduce the stored branches exactly. The largest term is
$0.62\sigma$ (FHC, open highest $p_\mu$ bin); for the derived observables it is at most $0.37\sigma$
(lower $W_\mathrm{had}$ region, combined). $\delta\phi_T$ and the observables that do not use the muon
momentum carry exactly zero. Released as \texttt{cov\_MCSscale.txt} and \texttt{cov\_total\_plusMCS.txt}
under \texttt{data\_release/cov/1p\_*} and summarised in \texttt{data\_release/mcs\_scale\_1p\_2026-09-27.tsv}
(\texttt{report/tools/mcs\_eval.py all 1p}); not yet inside the framework covariance.

\section{Document revision of 2026-09-27}

The three documents were revised for order and consistency. No result moved; the items below are
statements that were wrong or inconsistent and are now corrected, and material that moved between
documents.

\subsection{Beam-on control-region material held out}
The three documents are written as they stand before the control regions are compared with beam-on
data. The eighteen passages that used beam-on data (the comparison subsection of the analysis note
with its figure, and the sentences of the abstract, introduction, limitations, unblinding
requirements and summary, one clause of the proton-tagged note, and the control-region outcome and
one cross-check clause of the supplement) are kept verbatim, with their replacements, in
\texttt{report/beam\_on\_held\_out.tex} and \texttt{beam\_on\_held\_out.json} for merging back.
The analysis note label \texttt{sec:cr\_outcome} is replaced by \texttt{sec:cr\_procedure}.

\subsection{Inclusive analysis note}
\begin{itemize}\itemsep1pt
  \item The RHC paragraph of the results stated that no RHC dirt sample exists and that the FHC dirt
    overlay was scaled to the RHC exposure. The RHC configuration uses the RHC dirt overlays of
    Runs~3b, 4a and~4b, with Run~3b standing in for Runs~1 and~2; the text and the sample table now
    say so.
  \item The summary quoted $30$--$46\%$ bin-averaged uncertainty without data statistics and $67\%$
    for FHC $p_\pi$; the tables give $33$--$46\%$ and $57\%$. The column ``Total'' of the per-configuration
    result tables is the systematic total without the data statistics and is now labelled so.
  \item The detector appendix stated that track momentum is obtained from curvature (there is no
    magnetic field), that the NuMI spill is $1.6~\mu$s (that is the Booster spill; NuMI is about
    $10~\mu$s), that the signals are digitised inside the cryostat, and listed a cosmic BDT that the
    selection does not use. It was rewritten. The proton-tagged note no longer repeats it.
  \item Proton-tagged numbers were removed from the inclusive note (the proton-tagged column of the
    systematics table, the proton-tagged rows of the total cross-section table, the proton-tagged
    binning, cut parameters, alternative binnings and $W_{\pi p}$ appendix) and now appear only in the
    proton-tagged note, which has its own totals, binning and inventory tables. The generators
    (\texttt{rebuild\_tables.py}, \texttt{binning\_table.py}, \texttt{inventory.py}) write one table per
    note.
  \item The validation of the background fit (toys, injections, Markov chain, pull-mean convergence)
    moved to the control-region part of the technical supplement.
  \item Release dates were removed from the running text of the ensembles and appendices.
\end{itemize}

\subsection{Proton-tagged note}
\begin{itemize}\itemsep1pt
  \item Restructured in the order introduction, selection, observables and binning, extraction and
    uncertainties, results, validation, prerequisites.
  \item The text on the pair $\theta_{\pi p}\times W_{\pi p}$ quoted a conditioning of $0.560$, a
    $58\%$ data-statistical term, per-bin uncertainties up to $545\%$ and a closure ratio of $-0.40$,
    values of an earlier extraction; it now quotes the table ($0.537$, $78.3\%$, up to $276\%$, down
    to $-0.80$).
  \item The efficiency is quoted as $11.5\%$ (FHC), as in the cut-and-count table, not $11.4\%$; the
    largest detector term as $75\%$, not $100\%$.
\end{itemize}

\subsection{Technical supplement}
\begin{itemize}\itemsep1pt
  \item Reordered into parts on pion momentum and binning, proton-tagged observables, extraction and
    validation, control regions, and selection and infrastructure: the pion-estimator survey and the
    five-bin $p_\pi$ response moved to the pion-momentum section, the D'Agostini cross-check, the total
    cross section, the $A_C$ row sums and the ensembles to the validation part.
  \item The abstract stated that $p_\pi$ is reduced to two bins; it is measured in three regions.
  \item A reference intended for the extraction ensembles pointed to the background-fit toys.
  \item Prose revised throughout; tables, figures and numbers unchanged.
\end{itemize}

\section{Validation of the results proposed for approval (2026-09-28)}

Prerequisite V1 is complete. The results proposed for approval on 2026-09-27 that had not been through
the validation of the other extractions, $\theta_p$ and $\theta_{\pi p}$ (FHC, RHC, combined) and the
two combined-configuration two-dimensional results, $\theta_p\times\delta p_T$ and
$\cos\theta_\pi\times\cos\theta_\mu$, now have both tests.

\begin{itemize}\itemsep1pt
  \item \textbf{Alternative models.} GENIE multisim universe~545 and the $\Delta\to N\pi$ angular unisim,
    as for the other extractions: the proton angles in FHC and $\cos\theta_\pi\times\cos\theta_\mu$
    (SLURM 3433849), $\theta_p\times\delta p_T$ (SLURM 3433979), in the combined configuration for the
    two-dimensional results (\texttt{xsec\_analyzer/run\_altmodel\_v1.sh}, \texttt{configs/alt},
    \texttt{report/tools/altmodel\_eval.py V1}). No bin moves by more than $0.79\sigma$, the second
    $\theta_p$ bin under the GENIE variation, which returns all four $\theta_p$ bins $9$--$28\%$ low while
    their truth moves by $-3\%$; the two-dimensional results stay within $0.74\sigma$ (proton-tagged) and
    $0.61\sigma$ (inclusive).
  \item \textbf{Ensembles.} $100$ members for each of the eight extractions (SLURM 3433855--3433862,
    \texttt{xsec\_analyzer/ens\_v1.sh}): statistical pull widths $0.94$--$1.04$, coverage nominal, no
    truncated read in the $4800$ inputs ($33$ inputs of the inclusive two-dimensional ensemble are one
    entry short, the boundary case of an event at $\cos\theta_\mu=1$ exactly). In the sparsest inclusive two-dimensional cell
    ($\cos\theta_\mu<0.65$, $\cos\theta_\pi<0.35$) the ensemble mean lies $11\%$ below the reference,
    $0.58$ statistical and $0.17$ total standard deviations, and the ensemble-mean integral lies $1.2\%$
    below it ($3.4$ standard errors of the mean); no correction is applied
    (\texttt{data\_release/ensemble\_v1\_2026-09-28.tsv}).
  \item \textbf{Status.} V1 is removed from the prerequisites of the three products in
    \texttt{report/status.tsv}, and the prerequisite table marks it complete.
  \item \textbf{Tools.} The statistical-covariance step (\texttt{slurm/ens\_statcov.sh}) and the pull
    evaluation (\texttt{report/tools/ensemble\_stat\_pulls.py}) accept two-dimensional extractions, whose
    pulls are read from the all-slice closure file. The proton-tagged ensemble table has one row per
    extraction and includes the new ones (\texttt{report/tools/ensemble\_tables\_v1.py},
    \texttt{ensemble\_tables\_1p.py}).
\end{itemize}

\subsection{Documents}
\begin{itemize}\itemsep1pt
  \item The authors' revision of the wording of the three documents is adopted; no number changed.
  \item The label ``Secondary result proposed for approval'' was made plural by appending ``s'', which
    printed ``\ldots\ proposed for approvals''; a plural form is now defined in
    \texttt{shared/status\_vocab.tex}.
  \item Four tables that ran into the margin were narrowed.
\end{itemize}

\section{Approval material moved to a separate document (2026-09-28)}

The analysis note, the proton-tagged note and the technical supplement describe the analysis as it
stands. The material of the approval process moved to a new document, \texttt{approval\_request.tex}:
\begin{itemize}\itemsep1pt
  \item the decision requested from the collaboration (no longer boxed), the status tables with their
    prerequisites column, the prerequisite table with the criteria of U1--U7, P1--P4 and R1, and the
    validation ladder;
  \item from the supplement, the request to inspect the control regions, the frozen analysis state, the
    points the collaboration is asked to ratify, the approval checklist, the approval rules within the
    frozen procedure and the sign-off.
\end{itemize}
In the three documents the validation ladder is replaced by paragraphs of text, the sections on
prerequisites by sections on limitations, and statements of what is still needed by statements of what
has and has not been done. The analysis note starts with the introduction, followed by the measurement
definition; its status table is in the results section. The status tables of the three documents no
longer carry a prerequisites column (\texttt{report/tools/inventory.py} writes both versions; the closing
quotes of ``yes'' in their captions were lost to a quoting error and are restored). The three boxed
statements of the supplement are text. \texttt{report/tools/check\_documents.py} now refuses a
prerequisite identifier in an analysis document and accepts the decision list, the prerequisites and
the validation ladder only in \texttt{approval\_request.tex}. The status label ``Secondary result
proposed for approval'' is shortened to ``Secondary result'' in the documents, the release index and the
checks. No number changed.

\section{Collaboration review of 2026-09-28: corrections and shorter documents}

The review was answered without adding material; the three documents are shorter (analysis note $63\to61$,
proton-tagged note $22\to21$, supplement $78\to56$ pages). No result changed.
\begin{itemize}\itemsep1pt
  \item \textbf{Corrections.} The proton-tagged flux arithmetic now states that $298/331$ is the
    flux-dependent neutrino background over the signal of the total-cross-section table (whose other $7$
    background events are beam-off and dirt), and that the propagated universes give the remaining
    $1.9\%$ of the $35.5\%$. The proton-tagged combined integrals lie between FHC and RHC for nine of eleven
    extractions (not ten of twelve), and span $0.38$--$0.57$.
  \item \textbf{Wording.} The abstract of the analysis note puts the pseudo-data status first and no
    longer says ``we measure''; the reweighting test is called a test of reweighted interaction-model
    variations, and the limitations state that it cannot change final-state multiplicities; the
    proton-tagged total is described as the cross section of the final states with a proton above
    $0.3\GeVc$, and the generator spread as about one standard deviation of the measured total; ``response
    matrix'' replaces ``smearceptance'' in the prose. Total-cross-section captions state that the $60$~MeV
    flux includes the flux below the pion-production threshold; validation captions state that biases
    and full pulls use the framework covariance without the MCS term; the four statuses are defined in
    one sentence (\texttt{\textbackslash statDefinitions}).
  \item \textbf{Text extraction.} Every preamble maps glyphs, ligatures included, to Unicode
    (\texttt{glyphtounicode}), so copied text no longer contains ligature characters.
  \item \textbf{Shorter.} Supplement: the physics-background part removed (the analysis note keeps its
    physics context), the rejected alternatives condensed to one record each, the control-region machinery
    condensed (the stale FHC yields of the former uncertainty table and of the $\pi^0$-cut study, open item
    D1, went with them). Analysis note: the detector appendix in two paragraphs, the repeated limitation
    paragraphs and eight uncited references removed. Proton-tagged note: one results table for the three
    configurations (\texttt{rebuild\_tables.py}, \texttt{check\_tables.py}).
\end{itemize}

\section{Version 1.6 for the Editorial Board (2026-09-29)}

The authors' revision of the three documents is adopted and tagged as version~1.6 (git tag
\texttt{v1.6}). The beam-off normalisation is described by trigger rather than gate counts; the
status tables, the status labels of the proton-tagged note, the documents diagram, the glossaries of the
notes and the subsection on the comparison with beam-on data are removed; the analysis note states the
publication strategy that the separation into documents reflects. The supplement has its own preamble;
a duplicated block of cross-reference commands in it was removed, and references to the removed tables,
appendix and subsection were redirected or dropped. No number changed.

\section{Dirt normalisation in the count-level tables (2026-10-01)}

The cut-flow and cut-and-count macros (\texttt{cutflow\_yields.C}, \texttt{cutflow\_yields\_1p.C},
\texttt{total\_xsec\_counting.C}) scaled the Run-1 dirt sample by the data POT over the summed overlay POT
($0.081020$ FHC, $0.071666$ RHC), a placeholder introduced on 2026-08-09, instead of over the POT of the dirt
samples; the last two also applied the $0.65$ dirt weight twice. The macros now scale each run's dirt
sample by the data POT over its own POT, as the extraction does. The extraction is unchanged: the released
one-bin FHC total already subtracts a background of $590.0$ events, $6.6$ of them dirt. The dirt at the final
cut becomes $6.6$, $9.5$ and $16.0$ events (FHC, RHC, combined) instead of $1.1$, $1.0$ and $2.1$; the
purities become $58.5$, $57.4$ and $57.9\%$ ($51.9$, $49.7$ and $50.7\%$ proton-tagged). The cut-flow tables,
the cut-and-count tables, the cut-flow figures and the quoted purities are regenerated. The control-region
macros carried the same placeholder; their correction follows.

\section{Dirt normalisation in the control regions (2026-10-01)}

The control-region macros (\texttt{sideband\_compare.C}, \texttt{sb\_protocol.C}, \texttt{sb\_protocol\_data.C},
\texttt{sb\_plots.C}, \texttt{sb\_plots\_data.C}, \texttt{sb\_perrun\_data.C}, \texttt{bkgfit\_templates.C}) and three
scripts found with them (\texttt{bkgfit/farsb.py}, \texttt{bkgfit/run1\_trigger.py},
\texttt{pi0final/sb\_pi0\_final\_yields.C}) used the same placeholder, the $0.65$ applied twice in most of them.
All now take each run's dirt sample scaled by the data POT over its own POT, from the file list of the release
(\texttt{macros/dirt\_perrun.h}, \texttt{bkgfit/dirt.py}). The control-region dirt rises seven- to twentyfold:
$342/460/803$ events in CC$0\pi$ (FHC/RHC/combined), $93/124/216$ in $\pi^0$, $18.5/35.8/54$ in the cosmic region
and $0.5/1.7/2.1$ in multi-$\pi$. The neutrino and beam-off components, the selection and the cross sections are
unchanged. In the analysis note the dirt column of the control-region yields, the purity of the cosmic region
($4.1\%$) and of the selection without the opening-angle cut ($54.2\%$), the beam-off share of the cosmic region
($33/26\%$, now of the complete prediction) and its statistics ($283$ events, $21\%$) change. In the supplement the
full-stack test ($0.96$--$1.12$), the frozen-procedure statistics on pseudo-data (global $\chi^2=2.29/8$, pulls
within $\pm0.8$, shape tests $p>0.12$, single-region power $12\%$ and $44$--$62\%$, common-shift blind spot $50\%$,
null calibration $5.2\%$ and $1.1\%$), the shares of the fixed components, now with the dirt, and the
background-fit validation change. These numbers, the proton-tagged control-region table and the beam-off shares had
not been regenerated after the Run-1 exposure correction of 2026-09-19 and now include it. With beam-on data
(control-sample note) the frozen global test gives $\chi^2=6.03/8$ ($p=0.64$; it was $8.11/8$), the joint fit
gives CC$0\pi$ at $0.975\pm0.191$, the FHC Run-1 residual is unchanged ($0.862\pm0.021$), the cosmic region
moves from $1.21$ to $0.98$ and its far versions fall to $0.79$ and $0.76$. The supplement's comparison with the
$\nu_e$ analysis listed a $100\%$ dirt normalisation uncertainty among the shared conventions; this analysis
carries none, and the comparison now says so, with the bound that a $100\%$ variation moves the one-bin totals
by at most $0.9\%$.

\section{Version 1.6.1 for the Editorial Board (2026-10-01)}

The three documents with the two dirt corrections above are tagged as version~1.6.1 (git tag
\texttt{v1.6.1}), an errata release of version~1.6. The cross sections, the covariances and the data release
are those of version~1.6; the release notes list every changed number.

\section{Bragg-pion requirement dropped from the pion identification (2026-10-02)}

The pion identification of the single-pion selections required a Bragg-pion score of at least $0.08$. The
beam-on, beam-off and dirt ntuples of the older production (FHC Run~1, RHC Runs~1 and~3b) do not carry this
score, so the requirement passed every track in those samples while every overlay applied it. In the opened
control regions the released identification passed $7$--$11\%$ more beam-on candidates than predicted in those
three periods and not in the others; without the requirement the periods agree, and where the data carry the
score the requirement, after the other cuts, is modelled to within $2.2\sigma$. At unblinding the one-bin cross
sections would have been $4.7\%$ (combined, inclusive) and $4.6\%$ (proton-tagged) high. The requirement is
dropped in every sample (commit \texttt{1fb6be9}); the pion BDT, LLR and MIP-BDT requirements are unchanged.

The selection keeps $18.3\%$ of the signal at $57.3\%$ purity in FHC ($18.4\%$ and $56.6\%$ in RHC, $18.4\%$ and
$56.9\%$ combined; previously $17.5$--$17.6\%$ at $57.4$--$58.5\%$), and the proton-tagged selection $12.1\%$ at
$50$--$52\%$ (previously $11.5\%$). The pion candidate is a proton in $21\%$ of the selected events ($18\%$), and
the CC$0\pi$ background rises from $215$ to $260$ events in FHC, $42\%$ of the neutrino background. The one-bin
totals become $1.042$, $0.925$ and $0.975$ (FHC, RHC, combined; previously $1.034$, $0.932$, $0.975$), the
differential integrals move by factors of $0.92$--$1.12$ (median $1.00$), and the mean inclusive closure ratio is
$1.003$.

Every overlay, beam-off, dirt and detector-variation sample of both trees was reprocessed, the pseudo-data were
re-thrown with the release seeds (the same events; only the selection flags change), the $65$ universe files were
rebuilt with the bin definitions read back from the previous files, and the release chain, the D'Agostini
cross-check, the cut-flows, the tables and the figures were regenerated. The proton-rejection figure now shows the
pion-BDT score. The final-cut topology figure, which still showed a NuWro pseudo-data set at the superseded Run-1
exposure, is redrawn from the release samples. The counting cross-check now takes its systematic fractions from
the current $p_\mu$ results, and the proton-tagged cut-flow dirt carries the proton-tagged counters (it had been
processed with the inclusive selection). The binnings still meet the $0.50$ migration criterion (lowest diagonal
$0.516$, the inclusive three-region $p_\pi$ in FHC).

Carried over from version~1.6.1, made with the previous selection: the ensembles, the study extractions and the
binning-review studies. The MCS momentum-scale term and the alternative-model closure are redone (next section).
Dropping the requirement moves events from the control regions into the signal region: about $6$--$7\%$ of the new
signal region of the Run-4 and Run-5 beam-on data (about $85$ predicted events) was in the opened control regions.
The beam-on control-region skims made with the previous selection are withdrawn from use and have been remade; in
Runs 1--3 the beam-on selection does not change.

\section{Muon momentum-scale term and alternative-model closure with the new selection (2026-10-03)}

The data-side MCS term was evaluated again on pseudo-data re-thrown from the reprocessed samples with the release
seeds. In $p_\mu$ it is at most $0.60\sigma$ (FHC, open highest bin; previously $0.56\sigma$), $0.44\sigma$ combined
and $0.32\sigma$ in RHC; for FHC $p_\mu$ it is $9.7\%$ bin-averaged (previously $8.2\%$) and the total uncertainty
$41.4\%$. For the proton-tagged observables it is at most $0.66\sigma$ ($p_\mu$, FHC), and $0.38\sigma$ for the
derived observables. The official covariances, the summary tables and the systematic fractions of the counting
cross-check are regenerated. The statement that the downward scale leaves the two-bin $\delta p_T$ unchanged no
longer holds and is removed.

The alternative-model closure was repeated with the same universes and seeds. The largest inclusive bias is
$1.10\sigma$, in the third FHC $\cos\theta_\mu$ bin under the multisim variation (previously $1.18\sigma$ in the same
bin under the angular variation), and RHC stays within $0.85\sigma$. The low-imbalance proton-tagged $\delta p_T$ and
$p_n$ bins move by $-0.99\sigma$ and $-0.97\sigma$ under the multisim variation (previously $-1.14\sigma$ and
$-1.16\sigma$); the proton angles stay within $0.79\sigma$, and the two-dimensional results within $0.74\sigma$
(proton-tagged) and $0.55\sigma$ (inclusive). The note stated that the third FHC $\cos\theta_\mu$ bin has the
smallest $A_C$ row sum of the eight; the second bin has, in version~1.6.1 as now, and the statement is removed.

The truth-level changes of the two alternative models are now the ratio of the varied to the central-value
universe in the release universe files. They had been read from the truths of two independent throws, which adds
several percent of Poisson noise per bin. Under the multisim variation the signal rate changes by $3.5\%$ (FHC)
and $2.5\%$ (RHC) and the released $\cos\theta_\pi$ and $p_\pi$ bins by up to $30\%$ and $22\%$ (previously quoted as
$4\%$, $29\%$ and $24\%$); the RHC $p_\mu$ truth rises by $2$--$3\%$ (previously $3$--$9\%$), the near-threshold
FHC $p_\pi$ truth falls by $18$--$20\%$, the true cells of $\cos\theta_\pi\times\cos\theta_\mu$ change by $-23\%$ to
$+16\%$, and the proton-tagged low $\delta p_T$ and $p_n$ truths rise by $2\%$.

\section{Report folder organised by kind (2026-10-04)}

The analysis documents, their shared inputs (\texttt{shared/}, \texttt{glossary.tex}, \texttt{common\_preamble.tex},
\texttt{status.tsv}) and the control-region note are in \texttt{report/notes/}; papers, slides, other notes, plans,
result tables, figure macros, release packages and earlier exports are in \texttt{papers/}, \texttt{slides/},
\texttt{other\_notes/}, \texttt{planning/}, \texttt{results/}, \texttt{macros/}, \texttt{releases/} and \texttt{archive/}, and
\texttt{report/README.md} lists them. Figures, tables, the data release and the tools are where they were. Every
release table was regenerated in the new layout: the inclusive and proton-tagged $p_\pi$ region tables of the
supplement had not been regenerated for version~1.7 and now carry its values, as the corresponding figures did.

\section{Version 1.7 (2026-10-03)}

The three documents with the changes above are version~1.7. The cross sections, covariances and data release are
regenerated; the release notes list the changed numbers.

\section{Open items}

The following remain outstanding following the 2026-09-07 review. None is a correction awaiting application;
they are prerequisites, limits of the exposure, or tests not yet made. The analysis note's
request checklist carries the gating items; its Limitations section carries the tests not
yet made.

\begin{itemize}\itemsep2pt
  \item \textbf{Before the control regions are opened:} ratification of the frozen
    procedure by the conveners, including the $\pi^0$ definition and target class
    (adopted by the analysers 2026-09-06) in the light of the multiplicity-cut comparison
    now available; independent verification of the
    control-region skim (implemented 2026-09-06) by a second person, and the hashes of
    the frozen skim, plotting list and protocol macro in the approval record (repository
    tag \texttt{cr-freeze-2026-09-06}); the decision whether the $\pi^0$ region is
    diagnostic-only or eligible for a fitted correction, taken after the pre-fit result.
  \item \textbf{Before the signal region is opened:} confirmation that the flux constants
    are the PPFX-corrected central values; a NuWro or GiBUU fake-data pass through the
    released configurations (no alternative-generator test exists on the released chain);
    the clean-cache bit-identical reproduction repeated on the 2026-09-06 release; the inverted-proton-PID
    control region for any proton-tagged result; the documented treatment of any
    control-region discrepancy under the frozen procedure.
  \item \textbf{Ensembles:} the five rebinned extractions (FHC $p_\mu$, FHC and RHC
    $\cos\theta_\mu$, combined $p_\pi$, proton-tagged FHC $W_{\rm had}$) have their own $100$-member
    ensembles (2026-09-27): statistical pull widths $0.94$--$1.07$, coverage nominal, offsets
    $0.4$--$1.1\%$ at $\le2.0\sigma$; the one-bin totals and the proton-tagged $\delta p_T$ and $p_n$
    keep their binnings and their ensembles. The results proposed for approval on 2026-09-27 have
    theirs (2026-09-28). No test of this kind remains outstanding.
  \item \textbf{Systematic terms not in the released covariance:} the MCS muon momentum-scale
    term is evaluated at $\pm5\%$ on the data side and released as a separate component for $p_\mu$
    in all three configurations (up to $0.60\sigma$ per bin in FHC, $0.44\sigma$ combined and
    $0.32\sigma$ in RHC), is exactly zero for the other inclusive observables, and is not propagated
    inside the framework; the beam-off gate-ratio and Run-2 stand-in terms ($\lesssim0.4\%$
    of the total); the beamline-geometry flux term ($\lesssim0.2\%$).
  \item \textbf{Muon-momentum shape in the control regions:} the two-dimensional comparison with
    beam-on data shows a residual organised in $p_\mu$ (data over prediction $1.16$--$1.82$ at
    $0.35$--$0.75$~GeV/$c$) in both modes and in regions of different composition; open.
  \item \textbf{FHC Run~1 residual:} the $13\pm2\%$ offset of FHC Run~1 relative to the other five
    periods in the control regions, which survives the exposure correction and is not covered by any
    correlated term; its cause (bookkeeping POT, completeness of the Run-1 beam-on sample against its
    run list, or the Run-1 overlay normalisation) is not identified.
  \item \textbf{Limits of the exposure:} the differential $W_{\pi p}$ shape is withdrawn
    (no binning satisfies both the migration and the statistics requirement); the
    proton-tagged sample supports an integrated cross section.
  \item \textbf{Ensemble offset:} the ensemble means sit $0.4$--$1.1\%$ below the reference
    at up to $2.0\sigma$ (six ensembles, $100$ members each; the $2.0\sigma$ is the RHC total, $-0.8\%$); the throw itself is unbiased (mean thrown
    selected count $0.999$ of the expectation), so the offset, if real, is in the
    extraction; a larger ensemble would settle it.
\end{itemize}

\paragraph{External reproduction of the $A_C$-smeared comparison (resolved).}
The transformation
\[
  \left(d\sigma/dx\right)^{\rm smeared}_i = \frac{1}{w_i}\sum_j (A_C)_{ij}\,p_j ,
\]
applied to a prediction supplied as a per-bin cross section, reproduces every released
generator curve to a worst relative per-bin deviation of $4\times10^{-8}$, the precision of the
text files; it was the failure of
this test that exposed the double-$A_C$ bug, and its success is now the regression check.
No packaged reference implementation ships with the release; a user can verify their
implementation against any model column of \texttt{curves\_*.tsv}.

\paragraph{A caveat that is not an open item but reads like one.}
The closure $\chi^2$ values have $p$-values clustered near unity (median $0.98$ over the
fifty-one extractions) because the measurement and the truth it is compared against share
their systematics, so those terms cancel in the difference while still inflating the
covariance used to judge it. The ensemble of throws sizes the statistical part of the
band correctly; the systematic part remains untested by construction.

\section{Superseded values that may still be encountered}

\begin{center}\small\setlength{\tabcolsep}{4pt}
\begin{tabular}{p{4.3cm}p{4.2cm}p{6.4cm}}
\toprule
Quantity & superseded value & status \\
\midrule
Mean inclusive closure ratio        & $1.106$, $1.001$, $1.016$, $1.018$, $1.010$, $0.996$ & now $1.003$ (version~1.7) \\
Inclusive purity, final cut         & $58.5/57.4/57.9\%$ & $57.3/56.6/56.9\%$ without the Bragg-pion requirement (version~1.7) \\
Inclusive selection efficiency      & $17.5$--$17.6\%$ & $18.3$--$18.4\%$ (version~1.7) \\
Largest alternative-model bias      & $1.18\sigma$ (inclusive); $-1.14/-1.16\sigma$ (proton-tagged $\delta p_T$/$p_n$) & $1.10\sigma$; $-0.99/-0.97\sigma$ (version~1.7) \\
MCS term, FHC $p_\mu$              & $8.2\%$ bin-averaged, $0.56\sigma$ in the open bin & $9.7\%$, $0.60\sigma$ (version~1.7) \\
RHC $W_{\pi p}$ cross section       & $-0.194$         & sign was a POT-convention artefact \\
Combined detector systematic        & $42.3\%$         & $24.6\%$ at the fix; $21.4\%$ ($p_\mu$) in the current release \\
Combined result / expected          & $0.910$          & $1.027$ at the fix (2026-08-28 pass) \\
FHC $\theta_\mu$ $A_C$ conditioning & $0.015$          & now $0.51$ \\
Mean $\delta p_T$                   & $0.742$~GeV/$c$  & now $0.264$~GeV/$c$ (beam frame) \\
Proton-tagged \texttt{detVar}       & \texttt{n/a}, then $10.5$--$64.4\%$, then $10.5$--$108.5\%$ (2026-09-06) & now $9.9$--$74.5\%$ (median $30.9\%$, 2026-09-26 release) \\
Proton-tagged $\sigma$, all obs.    & 2026-08-18 values & superseded 2026-08-31, and again by the re-thrown release of 2026-09-04 (current values in the note) \\
Cut-and-count totals                & $1.06/0.98/1.02$; $1.03\pm0.43$ / $0.93\pm0.42$ / $0.98\pm0.40$ (borrowed systematic) & now the one-bin extraction $1.03\pm0.40$ / $0.93\pm0.35$ / $0.97\pm0.37$ \\
Proton-tagged purity                & $47$--$48\%$     & now $52/50/51\%$ \\
Signal-region dirt (final cut)      & $1.1/1.0/2.1$ (FHC/RHC/comb.) & now $6.6/9.5/16.0$; inclusive purity $58.8/57.7/58.2\%$ now $58.5/57.4/57.9\%$ (2026-10-01) \\
Control-region dirt (FHC)           & $49/0.1/10/1.9$ (CC$0\pi$/multi-$\pi$/$\pi^0$/cosmic) & now $342/0.5/93/18.5$ (2026-10-01) \\
Control-region global $\chi^2$      & pseudo-data $0.66/8$; beam-on $8.11/8$ & now $2.29/8$ and $6.03/8$ (2026-10-01) \\
Ensemble offset and pull widths     & $2.5\%$ at $2.1\sigma$; widths $1.04/0.96$ & now $0.5$--$2.7\%$ at $\le1.9\sigma$; widths $1.08/0.97/1.02$; 2026-09-20: six ensembles (FHC/RHC/COMB, $46$--$100$ members) $0.4$--$1.3\%$ at $\le1.5\sigma$, widths $0.90$--$1.10$; 2026-09-24: $100$ members each, race-corrupted members re-run, $0.4$--$1.1\%$ at $\le2.0\sigma$, widths $0.92$--$1.07$; 2026-09-27 (binnings of the 0.50 criterion): widths $0.95$--$1.07$ inclusive, $0.94$--$1.08$ proton-tagged, offsets $0.3$--$1.1\%$ at $\le2.0\sigma$ \\
FHC beam-off in the signal region   & $128$, then $13$, then $30$ & now $23$ (per-run gates, corrected Run-1 exposure) \\
Combined generator weights (FHC:RHC) & $0.458{:}0.542$ & now $0.425{:}0.575$ (2026-09-24) \\
Two-region $p_\pi$, FHC            & $0.059$ / $0.764$ & now $0.037$ / $0.827$ \\
Proton candidate $=$ leading proton  & $59\%$ (and $31\%$ ``not a proton'') & now $81.7\%$ ($12.6\%$ not a proton); the earlier figure predates the proton-tagged reprocess \\
$\theta_p$ bin edges (first)        & $0,\,0.597,\,0.942,\,1.287,\,\pi$ & replaced by RHC-derived edges (2026-09-18), then by the maximin ones (2026-09-26) \\
Control-region detector term        & $\sim\!20\%$ on CR totals & $\sim\!2\%$; the sideband \texttt{detVar} ratios were not POT-normalised \\
CR muon-angle shape failure         & reported in CC$0\pi$ and $\pi^0$ & withdrawn: the sideband MC predated the beam-frame correction \\
Finest $\cos\theta_\mu$, $\theta_\mu$ binning & five bins, ``cannot go finer'' & released five fail $0.68$; optimised edges admit more (2026-09-26 review) \\
$W_{\rm had}$ six-bin response       & ``adequate''     & fails $0.68$ in all six bins, $0.50$ in five; two regions admissible \\
Binning information measure         & $s_{\min}/s_1\gtrsim0.1$ & Wiener filter factors of $A_C$ and their sum; most released binnings fail the old threshold \\
Binnings (2026-09-26)               & $p_\mu$ 7, $\cos\theta_\mu$ 5, $p_\pi$ 2, $W_{\rm had}$ 6, $\delta\phi_T$ 2, $\theta_{\pi p}$ 3 bins & now 6, 8, 3 (inclusive), 2, 3, 4; $\theta_\mu$ dropped; $p_p$ added (3) \\
$\theta_p$ bin edges                & $0,\,0.659,\,1.068,\,1.539,\,\pi$ (RHC-derived) & now the maximin $0,\,0.44,\,0.785,\,1.162,\,\pi$~rad \\
Mean inclusive closure ratio        & $1.011$ (eighteen extractions) & $0.996$ (fifteen, 2026-09-26) \\
Smallest inclusive $A_C$ row sum    & $0.05$ (RHC $\cos\theta_\mu$, five bins) & $0.26$ (RHC $\cos\theta_\mu$, eight bins); now $0.33$ (RHC and combined $p_\mu$, version~1.7) \\
Identity $UR=A_C$                   & $4.4\times10^{-6}$ & $6.1\times10^{-5}$ (eight-bin $\cos\theta_\mu$, width ratio $56{:}1$) \\
Curve reproduction from $A_C$       & $1.2\times10^{-10}$ & relative $4\times10^{-8}$, the precision of the text files \\
\bottomrule
\end{tabular}
\end{center}

\end{document}
